% !TEX program = xelatex
% Compile twice with XeLaTeX.
\documentclass[11pt,a4paper]{article}

\usepackage{kotex}
\setmainhangulfont{Malgun Gothic}
\setsanshangulfont{Malgun Gothic}
\usepackage{amsmath,amssymb,mathtools,bm}
\usepackage{booktabs,longtable,tabularx,array}
\usepackage{geometry}
\geometry{margin=2.15cm}
\usepackage{caption}
\captionsetup{font=small,labelfont=bf}
\usepackage{enumitem}
\usepackage{xcolor}
\usepackage{fancyhdr}
\usepackage[hidelinks]{hyperref}
\hypersetup{
  pdftitle={Semianalytic Reconstruction of the Mean-Field Gain Model and the Spectral Fluctuation Theory for 85Rb Double-Lambda Four-Wave Mixing},
  pdfauthor={GABES validation report}
}

\definecolor{navy}{HTML}{183B56}
\definecolor{teal}{HTML}{207F79}
\definecolor{amber}{HTML}{A96500}
\definecolor{lightamber}{HTML}{FFF4DF}
\pagestyle{fancy}
\fancyhf{}
\lhead{\small $^{85}$Rb double-$\Lambda$ FWM}
\rhead{\small model audit}
\cfoot{\thepage}
\setlength{\headheight}{14pt}

\newcommand{\dd}{\mathrm d}
\newcommand{\ii}{\mathrm i}
\newcommand{\e}{\mathrm e}
\newcommand{\Tr}{\operatorname{Tr}}
\newcommand{\Rea}{\operatorname{Re}}
\newcommand{\Ima}{\operatorname{Im}}
\newcommand{\GHz}{\,\mathrm{GHz}}
\newcommand{\MHz}{\,\mathrm{MHz}}
\newcommand{\kHz}{\,\mathrm{kHz}}
\newcommand{\dB}{\,\mathrm{dB}}
\newcommand{\degC}{\,^{\circ}\mathrm C}
\newcommand{\chibar}{\bar\chi}
\newcommand{\ket}[1]{\lvert #1\rangle}
\newcommand{\bra}[1]{\langle #1\rvert}
\newcommand{\braket}[2]{\langle #1 \vert #2 \rangle}
\newcommand{\ketbra}[2]{\lvert #1\rangle\!\langle #2\rvert}
\newcommand{\pending}{\textcolor{amber}{not yet available}}
\newcommand{\pass}{\textcolor{teal}{PASS}}
\newcommand{\fail}{\textcolor{red!70!black}{FAIL}}
\newcolumntype{Y}{>{\raggedright\arraybackslash}X}
\newcolumntype{C}[1]{>{\centering\arraybackslash}p{#1}}
\newcolumntype{L}[1]{>{\raggedright\arraybackslash}p{#1}}
\setlist[itemize]{leftmargin=1.4em,itemsep=2pt,topsep=4pt}
\setlist[enumerate]{leftmargin=1.6em,itemsep=2pt,topsep=4pt}
\sloppy

\title{\textbf{Semianalytic Reconstruction of the Mean-Field Gain Model\\
and the Spectral Fluctuation Theory\\
for $^{85}$Rb Double-$\Lambda$ Four-Wave Mixing}\\[6pt]
\large Separation of semiclassical gain, quantum fluctuations, and external detection loss}
\author{GABES model-audit report}
\date{21 August 2026}

\begin{document}
\maketitle

\noindent
\fcolorbox{amber}{lightamber}{%
\begin{minipage}{0.94\textwidth}
\textbf{Scope and status.}
This document separates a semiclassical mean-field gain model, a
quantum-fluctuation propagation model, and an external detection model.
The mean-field chain and the quantum
fluctuation \emph{structure} are derived in closed form and verified
numerically. A separate slow reference now solves the self-consistent
pump-only state and infinitesimal trace-zero Nambu response for the standard
minus Raman branch. It is not substituted into the production finite-seed
gain path. A production frequency-dependent drift
$M_q(\Omega_{\rm SA})$ coupled to a microscopic atomic diffusion matrix
$D(\Omega_{\rm SA})$ is still not implemented, so no physically validated
squeezing spectrum is claimed here. Any squeezing number obtained from the
mean-field gains $G_p$ and $G_c$ alone is a diagnostic, not a prediction: the
associated two-mode map fails the canonical commutator test
($\mathcal G_{\rm flux}=G_p-(\omega_p/\omega_c)G_c\ne1$) in the dissipative
cell. Gains here follow the model's two-gain convention, in which $G_p$ and
$G_c$ are computed independently against the same probe-seed input; the
photon-flux gap, not the raw power-gain difference, is the canonical
commutator diagnostic.
The present claim gate is therefore: algebraic identities may pass,
quantitative apparatus-level gain is \textsc{unsupported}, and physical
squeezing is \textsc{unavailable} until the missing provenance, covariance,
and held-out validation gates are closed.
\end{minipage}}

\begin{abstract}
The report reconstructs the warm-vapor $^{85}$Rb D1 double-$\Lambda$
four-wave-mixing calculation in three explicitly separated layers:
(i) a semiclassical atomic-response and mean-field gain model,
(ii) a quantum-fluctuation propagation model with Langevin reservoirs, and
(iii) an external detection model. For the standard minus Raman branch, a
separate executable reference solves the pump-only state as a
finite-dimensional Liouvillian null-space problem and resolves weak
fluctuations at an independent analysis (sideband) frequency
$\Omega_{\rm SA}/2\pi=0$--$4\MHz$. The production gain path remains the
scan-converged finite-seed Floquet solve. The non-collinear two-photon Doppler
shift is derived and evaluated in a separate pump-only reference, but is not
yet included in either finite-seed production branch.
The physical susceptibility is dimensionless while the reduced response has
units of time; the Voigt prefactor, conjugate energy relation, and all $2\pi$
conventions are fixed explicitly.

Two analytic layers carry the frequency dependence. First, when the finite
Liouvillian is diagonalizable, the weak-field resolvent has an \emph{exact}
simple-pole representation: the atomic response is a sum of complex
Lorentzian terms whose isolated-pole centers are the \emph{negative} imaginary
parts and whose half-widths are the negative real parts of the Liouvillian
eigenvalues. A defective Liouvillian instead requires a Jordan or contour
representation and generally produces higher-order poles. Complete positivity
places an exact GKSL generator in the closed stable half plane while still
allowing peripheral modes. The numerical spectral check is therefore an
assembly invariant, not an independent validation of model quality. Second,
the distributed-noise covariance integral is solved as a differential
Lyapunov equation, whose solution is an entire
function of the Lyapunov superoperator. Thus its finite-$L$ coordinate-free
evaluation has no algebraic singularity at an exceptional point, although the
mode coalescence and parameter sensitivity of an exceptional point can still
be physical. Resolving the solution in the eigenbasis exhibits the same
trace--discriminant parametrization as the classical transfer matrix, while a
block-exponential evaluation avoids the ill-conditioned eigenbasis entirely.
Both results are verified against direct numerics at the $10^{-12}$ level or
better in their stated fixtures.

Applying the pole representation to the documented operating point locates the
slowest nonvanishing atomic pole at $0.520\MHz$ half-width at zero velocity,
about $5.1$ times the assumed ground-coherence rate, with per-class half-widths
of $0.22$--$1.49\MHz$ across the velocity distribution. The non-collinear
Doppler scale is instead a Gaussian rms width of about $1.38\MHz$; it cannot be
added in quadrature to those Lorentzian half-widths. A quantitative comparison
with the measured $3.5\MHz$ squeezing bandwidth still awaits
$S_-(\Omega_{\rm SA})$. A separate slow pump-only reference now performs the
two-dimensional Maxwell average: at $121\degC$ and $0.32^\circ$ it gives
$1.380174\MHz$ analytic rms and $1.380163\MHz$ by quadrature, while the
finite-seed production scan remains one-dimensional.

A Floquet audit finds that $N_F=1$ is not converged at the archived model
point: going to $N_F=2$ changes $G_p$ by $16.8\%$, changes the archived
conjugate transfer coefficient by $13.0\%$, and reverses the sign of the
corresponding historical coefficient gap; $N_F=2$ and $N_F=3$ agree.
The production core and compiled susceptibility kernel now implement the same
finite-order recurrence for arbitrary $N_F\geq1$. Reported seeded scans default
to $N_F=3$ and are accepted only after a full-scan comparison with $N_F=2$ of
the complex response and transfer coefficients, gains, wrapped phases, and the
probe-gain optimum.
The independent pump reference is gauge-equivalent to the minus-branch
pump-only Floquet state to $8.37\times10^{-15}$. Across a declared
$4\times5$ analysis-frequency/detuning scan on the full one-dimensional
Doppler grid, its maximum normalized response residual is
$7.32\times10^{-17}$. The finite-seed response approaches this reference
monotonically as the seed is reduced, but the fitted last-three-point
Rabi-amplitude slope is $1.749$ and is not reported as an exact quadratic law.
Its opt-in $(v_z,v_x)$ wrapper keeps the laboratory beat and analysis
frequency fixed, shifts only the atomic detunings, and passes the declared
grid and cutoff convergence gates.
Positivity of
the truncated state is shown \emph{not} to detect this error, so Floquet
convergence testing cannot be replaced by a physicality check. The literature
operating point $\Delta/2\pi\simeq+0.9\GHz$, $\delta/2\pi\simeq-8\MHz$, and
$T=121\degC$ is kept distinct from the archived reduced-model point
$(\Delta/2\pi,\delta/2\pi,T)=(-1.50\GHz,-280\MHz,110\degC)$. Agreement with
the production code is classified only as an algebraic and numerical
consistency check within shared assumptions.
\end{abstract}

\tableofcontents
\clearpage

%======================================================================
\part{Definitions and conventions}
%======================================================================

\section{Atomic references and frequency definitions}

The four hyperfine states are
\[
g_1=5S_{1/2}(F=2),\quad g_2=5S_{1/2}(F=3),\quad
e_2=5P_{1/2}(F'=2),\quad e_3=5P_{1/2}(F'=3).
\]
The optical reference is
\[
\omega_{13}\equiv(E_{e_3}-E_{g_1})/\hbar,\qquad
\omega_{\rm HF}\equiv(E_{g_2}-E_{g_1})/\hbar .
\]
Every detuning and linewidth in an equation is an \emph{angular} frequency.
Numerical values are always displayed after division by $2\pi$:
\[
\frac{\Delta}{2\pi}=-1.50\GHz,\qquad
\frac{\delta}{2\pi}=-280\MHz,\qquad
\frac{\Gamma}{2\pi}=5.746\MHz .
\]

The pump and probe-seed conventions are
\begin{align}
\Delta&=\omega_{\rm pump}-\omega_{13},\\
\omega_p&=\omega_{\rm pump}-\omega_{\rm HF}+\delta,\\
\Omega_{\rm beat}&=\omega_{\rm pump}-\omega_p=\omega_{\rm HF}-\delta,\\
\omega_c&=2\omega_{\rm pump}-\omega_p .
\label{eq:frequency-definitions}
\end{align}
Thus $\Omega_{\rm beat}$ is the optical pump--probe-seed beat that organizes the
Floquet harmonics. It is not the analysis frequency.

The independent radio-frequency analysis variable is
\[
\frac{\Omega_{\rm SA}}{2\pi}\simeq0.1\text{--}4\MHz .
\]
It labels quantum sidebands at $\omega_p\pm\Omega_{\rm SA}$ and
$\omega_c\mp\Omega_{\rm SA}$. With the Fourier convention used here,
\begin{equation}
\delta\hat a_j(z,t)=
\int_{-\infty}^{\infty}\frac{\dd\Omega_{\rm SA}}{2\pi}\,
\delta\hat a_j(z,\Omega_{\rm SA})\e^{-\ii\Omega_{\rm SA}t}.
\label{eq:fourier}
\end{equation}

\begin{table}[ht]
\centering
\caption{Frequency definitions used throughout the report.}
\label{tab:frequency-definitions}
\small
\begin{tabularx}{\textwidth}{L{2.5cm}L{4.2cm}L{3.3cm}Y}
\toprule
Symbol & Definition & Reported as & Role\\
\midrule
$\omega_{13}$ & $(E_{e_3}-E_{g_1})/\hbar$ &
$\omega_{13}/2\pi$ & atomic optical reference\\
$\Delta$ & $\omega_{\rm pump}-\omega_{13}$ &
$\Delta/2\pi$ & pump one-photon detuning\\
$\delta$ & $\omega_p-\omega_{\rm pump}+\omega_{\rm HF}$ &
$\delta/2\pi$ & Raman detuning\\
$\Omega_{\rm beat}$ & $\omega_{\rm pump}-\omega_p$ &
$\Omega_{\rm beat}/2\pi$ & optical pump--probe-seed beat\\
$\Omega_{\rm SA}$ & sideband offset from an optical carrier &
$\Omega_{\rm SA}/2\pi$ & RF analysis (sideband) frequency\\
$\omega_p,\omega_c$ & probe and conjugate carrier angular frequencies &
$\omega_{p,c}/2\pi$ & optical carriers with
$\omega_c=2\omega_{\rm pump}-\omega_p$\\
$-\Ima\lambda_m$ & isolated-pole center in the convention of
Eq.~\eqref{eq:fourier} & $-\Ima\lambda_m/2\pi$ & compared with
$\Omega_{\rm SA}/2\pi$, never with $\Omega_{\rm beat}/2\pi$\\
$\Gamma$ & population-decay angular rate &
$\Gamma/2\pi$ & natural linewidth\\
\bottomrule
\end{tabularx}
\end{table}

\paragraph{Mode labels.}
The three optical fields are the \emph{pump}, the \emph{probe} (Stokes), and
the \emph{conjugate} (anti-Stokes); the weak field injected on the probe leg
is the \emph{probe-seed}, and probe and conjugate together are the
\emph{twin beams} \cite{Sim2025}. Subscripts follow that naming:
$p$ denotes the probe and $c$ the conjugate, while pump quantities carry the
explicit subscript ``pump'' so that no symbol is overloaded. Powers are
$P_{\rm pump}$, $P_0$ for the injected probe-seed, and $P_p$, $P_c$ for the
two output beams. Macroscopic polarizations are written $\mathcal P_j$ to
keep them distinct from powers, and the quadratures of mode $j$ are
$(X_j,Y_j)$.

\section{Field, beam, and susceptibility normalization}

Canonical traveling-wave amplitudes are photon-flux normalized:
\begin{equation}
[\delta\hat a_j(z,t),\delta\hat a_k^\dagger(z,t')]
=\delta_{jk}\delta(t-t'),\qquad
P_j=\hbar\omega_j|\alpha_j|^2 .
\label{eq:flux-normalization}
\end{equation}
If the atomic code uses Rabi amplitudes
$\Omega_j=\mathcal G_j\alpha_j$, its propagation matrix must be transformed
before being interpreted quantum mechanically:
\begin{equation}
M_{\rm can}=\mathcal G^{-1}M_{\Omega}\mathcal G,\qquad
\mathcal G=\operatorname{diag}(\mathcal G_p,\mathcal G_c^*).
\label{eq:canonical-rescaling}
\end{equation}
The near equality of the two optical wavelengths does not remove this
normalization requirement. Every quantity below that is contracted with the
canonical metric $J=\operatorname{diag}(1,-1)$ --- in particular the
commutator-defect matrix of Sec.~\ref{sec:realizability} --- is
basis dependent through Eq.~\eqref{eq:canonical-rescaling} and must be
recomputed if the collected mode areas differ.

The quoted beam sizes
\[
w_{\rm pump}=530\,\mu{\rm m},\qquad w_p=330\,\mu{\rm m}
\]
are $1/\e^2$ \emph{intensity radii}, not diameters.

The reduced response is defined schematically by
$\chibar_{jk}=\rho_{eg}^{(j\leftarrow k)}/\Omega_k$. Consequently
\begin{equation}
[\chibar_{jk}]={\rm s},\qquad
\chi_{jk}=-\frac{2Nd^2}{\epsilon_0\hbar}\,\ell_{\rm norm}\chibar_{jk},
\qquad [\chi_{jk}]=1 .
\label{eq:susceptibility-dimension}
\end{equation}
The factor $\ell_{\rm norm}$ is dimensionless. Its decomposition is audited
in Appendix~\ref{app:normalization}. The structural factor is now
$1/[2(2I+1)]=1/12$: the trace-normalized $\rho_{\rm ss}$ supplies the
ground-manifold population, while the inherited $0.74$ remains an empirical
residual rather than a first-principles quantity.

\section{Operating-point classification}

\begin{table}[ht]
\centering
\caption{Operating points are classified by provenance, not by numerical
appearance in a scan.}
\label{tab:operating-points}
\small
\begin{tabularx}{\textwidth}{L{3.2cm}C{2.0cm}C{2.0cm}C{1.7cm}Y}
\toprule
Point and class & $\Delta/2\pi$ & $\delta/2\pi$ & $T$ & Meaning\\
\midrule
Archived reduced-model optimum &
$-1.50\GHz$ & $-280\MHz$ & $110\degC$ &
Numerical optimum of the gain-only squeezing objective. It is not an
experimental optimum and its $N_F=1$ result is unconverged.\\
Sim et al. literature benchmark \cite{Sim2025} &
$+0.9\GHz$ & $-8\MHz$ & $121\degC$ &
Experimentally measured operating point: $600\,{\rm mW}$ pump,
$8$--$10\,\mu{\rm W}$ probe-seed, $12.5\,{\rm mm}$ cell,
$w_{\rm pump}=530\,\mu{\rm m}$, $w_p=330\,\mu{\rm m}$, and
$\theta=0.32^\circ$.\\
\bottomrule
\end{tabularx}
\end{table}

The literature describes the gain as about $15$; the cited $^{85}$Rb squeezing
trace reports raw powers $(P_0,P_p,P_c)=(8,111,109)\,\mu{\rm W}$, which give
the power ratios $13.875$ and $13.625$ used here.
It reports maximum IDS of $-7.8\dB$, an approximate $3.5\MHz$
bandwidth, and a laser-source electronic-noise peak near $2\MHz$. Those
measurements are not properties of the archived
$(\Delta/2\pi,\delta/2\pi)=(-1.50\GHz,-280\MHz)$ model point.

%======================================================================
\part{Pump-only atomic steady state}
%======================================================================

\section{Self-consistent finite-dimensional Liouvillian}

For each velocity class and longitudinal position, the pump-only state is
defined by
\begin{equation}
\mathcal L_{\rm pump}(z,\bm v)\rho_{\rm ss}(z,\bm v)=0,\qquad
\Tr\rho_{\rm ss}(z,\bm v)=1 .
\label{eq:pump-nullspace}
\end{equation}
For the standard minus Raman branch, the seeded frame has
$\Omega_{\rm beat}=\omega_{\rm HF}-\delta$. The explicit transformation
\begin{equation}
U_-(t)=\exp\!\left(+\ii\Omega_{\rm beat}t\ketbra{g_2}{g_2}\right)
\label{eq:pump-frame-transform}
\end{equation}
is used with $\rho_{\rm pump}=U_-^\dagger\rho_{\rm seed}U_-$ and
$H_{\rm pump}=U_-^\dagger H_{\rm seed}U_- -\ii\hbar
U_-^\dagger\dot U_-$. It makes the second pump coupling static and changes the
$g_2$ diagonal from
$\delta$ to
$\delta+\Omega_{\rm beat}=\omega_{\rm HF}$. The resulting pump Hamiltonian is
\begin{equation}
\frac{H_{\rm pump}}{\hbar}=
\omega_{\rm HF}\ketbra{g_2}{g_2}
-\sum_e\Delta_e(\bm v)\ketbra{e}{e}
+\frac12\sum_{g,e}
\left(\Omega^{\rm pump}_{ge}\ketbra{e}{g}+\Omega^{{\rm pump}*}_{ge}\ketbra{g}{e}\right),
\label{eq:pump-hamiltonian}
\end{equation}
where $\Delta_e(\bm v)=\Delta_e-\bm k_{\rm pump}\cdot\bm v$.
Spontaneous emission, collisional optical dephasing, and ground-state
decoherence enter the same reference Liouvillian. The shared
dissipator now includes the trace-preserving thermal reload channel
$\gamma_t(\rho_{\rm th}\Tr\rho-\rho)$ with ground-manifold weights $5/12$ and
$7/12$; the density-dependent Rb--Rb term remains a separate pure-dephasing
channel. Vectorizing the
$4\times4$ density matrix gives a $16\times16$ null-space problem; one row is
replaced by the trace constraint. Solving that system makes populations
and coherences feed back on one another self-consistently.

The diagonal phase in Eq.~\eqref{eq:pump-frame-transform} leaves the declared
projector dephasing and thermal reload invariant; phases acquired by individual
collapse operators cancel in their Lindblad dissipators. The executable gauge
map reconstructs the seeded-frame pump state with $g_2$-row coherences in
$n=-1$, $g_2$-column coherences in $n=+1$, all remaining entries in $n=0$, and
all $|n|\geq2$ harmonics exactly zero. The reconstructed state agrees with the
direct $N_F=3$ pump-only Floquet solution to $8.37\times10^{-15}$. This
certification currently applies only to the standard minus Raman branch. The
inherited plus-branch finite-seed frame fails this gauge-parity test and the
pump-only reference therefore returns \texttt{unsupported} rather than mapping
it silently.

Equations~\eqref{eq:pump-nullspace}--\eqref{eq:pump-hamiltonian} are now an
executable \emph{semianalytic pump-only steady-state reference}: the
Liouvillian matrix elements are analytic functions of the parameters and the
null space is solved numerically. It remains separate from the current gain
tables. The production path instead solves a
trace-one, finite-seed density matrix at $N_F\geq3$ and applies an adjacent-order
full-scan convergence gate. The generalized audit separately repeats the same
finite-seed finite-Floquet construction for $N_F=1,2,3$.
The older rate/Sylvester carrier in
\path{analysis/squeezing/analytic_reconstruction/theory_final.py} is a third,
historical leading-order approximation; it is not the source of the corrected
Option-A tables. None of these distinctions changes the fact that any solve is
exact only within its four-level, Markov, rotating-wave, field, velocity, and
Floquet assumptions. In particular, the reference is undepleted, uses the
one-dimensional longitudinal Doppler grid by default, and does not recompute
the pump state segment by segment. Section~\ref{sec:doppler} adds a separate
two-dimensional wrapper around this reference; it does not promote that
wrapper into the finite-seed production path.

The finite-seed path takes the probe-seed power already present at the cell as
an input. It does not model the apparatus EOM $V_\pi$, insertion loss, RF-to-
sideband conversion, or sideband impurity. Those quantities belong to the
lab-to-cell input map and technical-noise budget, not to the pump-nullspace
equation.

\section{Validity of the historical rate/Sylvester factorization}

An alternative construction solves saturation rate equations for populations
and then solves a Sylvester equation for coherences. Because the resulting
coherences do not feed back into the populations, that construction is not an
exact strong-pump steady state. It may be used only under all of the
following assumptions:
\begin{itemize}
\item optical coherences can be adiabatically eliminated;
\item excited-state interference and coherent population trapping are small;
\item pump-induced level dressing is weak compared with relevant separations;
\item population redistribution is slow compared with optical dephasing; and
\item fitted rates represent the omitted coherent dynamics.
\end{itemize}
At the archived point the quoted pump Rabi rate is comparable to or larger
than the excited-state hyperfine splitting
($\Omega_{\rm pump}/2\pi=0.708\GHz$ versus
$\nu^{(e)}_{\rm HF}=361.58\MHz$), so these assumptions are not automatically
satisfied. All population and coherence
numbers obtained this way are therefore classified as a controlled
approximation only if a separate error study is supplied.

\section{Phenomenological rates and truncations}

The four-level pump model uses a Zeeman-averaged hyperfine basis. The
ground-coherence floor, density-dependent collisional term, and optical
self-broadening are phenomenological inputs unless independently measured.
At the literature point the engine evaluates
\begin{equation}
\frac{\gamma_{gg}}{2\pi}=102.882\kHz,\qquad
\frac{\gamma_{\rm opt}}{2\pi}=0.7416\MHz,\qquad
\frac{\Gamma}{2\pi}=5.7460\MHz,
\label{eq:rates-at-literature}
\end{equation}
for $N=2.1496\times10^{19}\,{\rm m^{-3}}$ at $T=121\degC$, where the effective
ground-coherence decay is the sum of the $100\kHz$ trace-preserving transit
reset and a $2.882\kHz$ Rb--Rb pure-dephasing term, and
$\gamma_{\rm opt}=\tfrac12\beta N$ is the added collisional optical half-width.
Thus the transit contribution replenishes ground populations rather than merely
damping the Raman coherence, while preserving trace exactly. Their provenance
and sensitivity are listed in
Appendix~\ref{app:parameters}. No calculation involving these quantities is
labeled exact.

%======================================================================
\part{Weak-field linear response}
%======================================================================

\section{Frequency-dependent response matrix}

Let $\mathcal V_k(\Omega_{\rm SA})$ be the perturbation generated by the weak
probe or conjugate sideband, and let $\omega_k(\delta,\Omega_{\rm SA})$ be that
field's angular frequency in the static pump frame. With
Eqs.~\eqref{eq:fourier} and
$\dot\rho=\mathcal L_{\rm pump}\rho+\mathcal V\rho_{\rm ss}$,
\begin{equation}
\delta\rho_k(\Omega_{\rm SA})=
-\left(\mathcal L_{\rm pump}+\ii\omega_k(\delta,\Omega_{\rm SA})\right)^{-1}_{\Tr=0}
\mathcal V_k(\Omega_{\rm SA})\rho_{\rm ss}.
\label{eq:linear-resolvent}
\end{equation}
For the implemented minus-branch Nambu sector,
\begin{equation}
\omega_k^{(-)}(\delta,\Omega_{\rm SA})
=-\omega_{\rm HF}+\delta+\Omega_{\rm SA}.
\label{eq:pump-frame-response-frequency}
\end{equation}
Thus the optical two-photon detuning and RF analysis frequency are independent
inputs; changing one never silently changes the other. The inverse acts on the
trace-zero subspace. At $\omega_k=0$ the implementation projects the full
stationary subspace instead of applying an ordinary singular inverse. This
zero relative resolvent frequency is not generally the same as setting the lab
analysis frequency $\Omega_{\rm SA}$ to zero. Using the opposite Fourier sign
would reverse the sign of the imaginary term; mixing those conventions is not
allowed.

Polarization readout gives the full response matrix
\begin{equation}
\begin{pmatrix}\mathcal P_p\\\mathcal P_c^*\end{pmatrix}
=\epsilon_0
\begin{pmatrix}
\chi_{pp}(\Omega_{\rm SA})&\chi_{pc}(\Omega_{\rm SA})\\
\chi_{cp}^*(\Omega_{\rm SA})&\chi_{cc}^*(\Omega_{\rm SA})
\end{pmatrix}
\begin{pmatrix}E_p\\E_c^*\end{pmatrix}.
\label{eq:chi-matrix}
\end{equation}
Thus $M_q$, the reservoir coupling $B$, and the diffusion matrix $D$ may all
depend on $\Omega_{\rm SA}$.
The executable reference returns the complete complex $2\times2$ derivative
with Nambu input/output order $(p,c^*)$ and units of seconds per unit angular
Rabi frequency. It supplies the atomic-response ingredient needed for a future
$M_q(\Omega_{\rm SA})$ construction, but it does not construct $B$ or $D$ and
is not the production gain solver.

\section{Non-collinear Doppler geometry}
\label{sec:doppler}

Choose the pump along $z$ and the probe-seed in the $x$--$z$ plane:
\[
\bm k_{\rm pump}=k_{\rm pump}\hat{\bm z},\qquad
\bm k_p=k_p(\sin\theta\,\hat{\bm x}+\cos\theta\,\hat{\bm z}).
\]
The laboratory detunings seen by an atom with velocity
$\bm v=(v_x,v_y,v_z)$ are therefore
\begin{align}
\Delta_{\rm eff}&=\Delta-k_{\rm pump}v_z,\\
\delta_{\rm eff}&=\delta+(k_{\rm pump}-k_p\cos\theta)v_z
-k_p\sin\theta\,v_x .
\label{eq:doppler-detunings}
\end{align}
The sign of the transverse term follows directly from
$\delta_{\rm eff}=\delta-(\bm k_p-\bm k_{\rm pump})\cdot\bm v$.
The quantitative average is at least two-dimensional:
\begin{equation}
\langle X\rangle_v=
\int\!\!\int\frac{\dd v_z\,\dd v_x}{2\pi\sigma_v^2}
\exp\!\left[-\frac{v_z^2+v_x^2}{2\sigma_v^2}\right]
X(\Delta_{\rm eff},\delta_{\rm eff}).
\label{eq:two-dimensional-velocity}
\end{equation}
The unused $v_y$ projection integrates to one. For
$k_{\rm pump}\simeq k_p=k$,
\begin{equation}
|\bm k_{\rm pump}-\bm k_p|=2k\sin(\theta/2)\simeq k\theta,\qquad
\frac{\sigma_{\delta,\theta}}{2\pi}
\simeq\frac{\theta\sigma_v}{\lambda}=1.361\MHz
\label{eq:angular-doppler}
\end{equation}
for $\theta=0.32^\circ$, $\sigma_v=193.7\,{\rm m/s}$, and
$\lambda=795\,{\rm nm}$. These values correspond to the $110\degC$ estimate. At the
$121\degC$ literature point, $\sigma_v\simeq196.5\,{\rm m/s}$ and the width is
about $1.38\MHz$. Both are much larger than
$\gamma_{gg}/2\pi=0.103\MHz$. The angular term vanishes as
$\theta\to0$; a roughly $2\kHz$ collinear residual remains because
$\omega_{\rm pump}\ne\omega_p$.

The minus-branch pump-only reference now implements
Eq.~\eqref{eq:two-dimensional-velocity} as a tensor Gauss--Legendre quadrature
of the truncated Maxwell distribution. Its interface keeps
$\Omega_{\rm beat,lab}$ and $\Omega_{\rm SA}$ independent of velocity and
applies Eq.~\eqref{eq:doppler-detunings} only to the atomic detunings. This
separation is essential: passing $\delta_{\rm eff}$ to the production Floquet
table would also move the laboratory beat frequency and would describe a
different problem.

At the literature point the analytic rms is $1.380173564\MHz$, of which
$1.380161171\MHz$ is transverse and $5.848714\kHz$ is axial at angle. The
$40\times40$, $5\sigma_v$ quadrature gives $1.380163304\MHz$, a
$-7.43\times10^{-4}\%$ difference. Table~\ref{tab:angular-doppler-convergence}
records the feature-level convergence. These large gains are mean-field
diagnostics under the same uncalibrated coupling assumptions as the other gain
tables, not experimental predictions.

\begin{table}[ht]
\centering
\caption{Two-dimensional non-collinear reference convergence. The refined
scan spacing is $0.05\MHz$; the first three rows use a $5\sigma_v$ cutoff.}
\label{tab:angular-doppler-convergence}
\small
\begin{tabular}{c r r r r}
\toprule
Grid or cutoff & velocity pairs & feature $\delta/2\pi$ & $G_p$ & change\\
\midrule
$24\times24$ & 576 & $36.15\MHz$ & 4517.4025 & $+0.01495\%$ vs 40\\
$32\times32$ & 1024 & $36.00\MHz$ & 4517.0481 & $+0.00711\%$ vs 40\\
$40\times40$ & 1600 & $35.95\MHz$ & 4516.7271 & reference\\
$4.5\sigma_v$ at 40 & 1600 & $35.95\MHz$ & 4519.0828 & $+0.05215\%$\\
$5.0\sigma_v$ at 40 & 1600 & $35.95\MHz$ & 4516.7271 & reference\\
\bottomrule
\end{tabular}
\end{table}

The accepted successive-grid comparison is $32\to40$: the gain changes by
$0.00711\%$ and the sampled feature position by $0.05\MHz$. Extending the
cutoff from $4.5\sigma_v$ to $5\sigma_v$ changes the gain by $0.05215\%$ and
does not move the sampled feature. At the declared $-8\MHz$ literature point,
the one-dimensional and two-dimensional diagnostics are respectively
$G_p=403.825150$ and $404.304420$, a $+0.118683\%$ change. The response residual
is $2.18\times10^{-16}$ and the trace error is $5.83\times10^{-26}$.

This closes the slow-reference quadrature task, not the production geometry:
the finite-seed scan remains one-dimensional, the crossing angle is a single
value rather than a beam angular distribution, and the pump state is not
recomputed segment by segment. No microscopic noise is added.

\section{Floquet harmonics and convergence requirement}

If the chosen rotating frame leaves a periodic pump term at
$\Omega_{\rm beat}$, expand
\[
\rho(t)=\sum_{n=-N_F}^{N_F}\rho_n\e^{-\ii n\Omega_{\rm beat}t}.
\]
The finite block equations are
\begin{equation}
\left(\mathcal L_0+\ii n\Omega_{\rm beat}\right)\rho_n
+\mathcal C_p\rho_{n-1}+\mathcal C_m\rho_{n+1}=0 .
\label{eq:floquet-block}
\end{equation}
Trace one is imposed in the $n=0$ sector and trace zero in every other sector.
Before truncation, the same blocks define an infinite-dimensional
Floquet--Liouville operator $\mathbb L$; its eigenvalues are Floquet exponents,
defined modulo $\ii\Omega_{\rm beat}$. When the relevant spectral decomposition
exists, the resolvent identities of Sec.~\ref{sec:pole-expansion} then apply on
that extended space, with every harmonic sector contributing a frequency
replica; a simple-pole sum additionally requires diagonalizability.

For the exact periodic GKSL problem, the one-period propagator is CPTP, so each
Floquet multiplier $r_m=\exp(\lambda_mT)$ obeys $|r_m|\leq1$. A finite-$N_F$
projected block matrix is not generally the generator or logarithm of a CPTP
map, however; positive real parts may be truncation artifacts, and passing a
half-plane or state-positivity test does not establish convergence. The
truncation is accepted only after convergence with respect to $N_F$ has been
established for the transfer coefficients, their wrapped phases, and the
specified optimum $\delta^*$.
The production gate uses the dimensionless mixed tolerance
$|\Delta T|\leq10^{-10}+0.01\max(|T_N|,|T_{N-1}|)$ for every transfer
coefficient. Reduced susceptibilities have units of time, so their absolute
floor is instead $10^{-8}$ times that component's full-scan maximum; this is
added to the same one-percent local relative term. Derived gains must change by
less than one percent, retained wrapped phases by less than $0.5^\circ$, and the
probe-gain optimum by at most one final-grid interval. A failed condition returns
\texttt{UNCONVERGED}.
The production core and compiled kernel now use generalized block continued
fractions and are pinned for $N_F=1,2,3$ against a separately assembled dense
finite block, including trace and normalized-equation residual tests. The
archived part of the audit in Sec.~\ref{sec:floquet-convergence} retains its own
assembly, reproduces \path{ref_solver.py}'s historical $N_F=1$ finite-seed
solve, and agrees with direct finite-block solves to $5.4\times10^{-13}$. Those
numbers isolate truncation under the historical dressed-$k$ convention; they are not production
Option-A physical predictions. The separate literature-point section uses the
current density-matrix assembly and checks only its Maxwell/$Q$ propagation
against a literal independent SI reconstruction. The separate pump-only
reference implements Eq.~\eqref{eq:pump-nullspace} for the minus branch and is
checked against both its gauge-mapped finite-Floquet state and an independent
dense/direct response construction; it is not substituted into either of those
finite-seed paths.

%----------------------------------------------------------------------
\section{Exact spectral representation of the resolvent}
\label{sec:pole-expansion}

Equation~\eqref{eq:linear-resolvent} is usually evaluated by a numerical
matrix inversion at each $\Omega_{\rm SA}$, which hides the analytic
structure of the response. The representation below exposes it without
approximation.

\subsection{Eigenbasis of the Liouvillian}

Let $\mathcal L\equiv\mathcal L_{\rm pump}$ act on the $n^2$-dimensional
vectorized space and let its right and left eigenvectors be
\begin{equation}
\mathcal L\,\ket{r_m}=\lambda_m\ket{r_m},\qquad
\bra{l_m}\mathcal L=\lambda_m\bra{l_m},\qquad
\braket{l_m}{r_{m'}}=\delta_{mm'} ,
\label{eq:biorthogonal}
\end{equation}
where $\braket{\cdot}{\cdot}$ is the Hilbert--Schmidt pairing. $\mathcal L$ is
not Hermitian, so left and right eigenvectors differ and the eigenvalues are
complex. Whenever $\mathcal L$ is diagonalizable, the resolvent has the exact
partial-fraction form
\begin{equation}
\boxed{\;
\left(\mathcal L+\ii\Omega\right)^{-1}
=\sum_{m}\frac{\ketbra{r_m}{l_m}}{\lambda_m+\ii\Omega}\; .}
\label{eq:resolvent-spectral}
\end{equation}
Substituting into Eq.~\eqref{eq:linear-resolvent} and contracting with the
Clebsch--Gordan-weighted polarization readout $\bra{w_j}$ gives the reduced
response as a pole sum,
\begin{equation}
\boxed{\;
\chibar_{jk}(\Omega)=\sum_m\frac{A^{(jk)}_m}{\lambda_m+\ii\Omega},
\qquad
A^{(jk)}_m=-\braket{w_j}{r_m}\bra{l_m}\mathcal V_k\ket{\rho_{\rm ss}} .}
\label{eq:chi-pole-sum}
\end{equation}

Equation~\eqref{eq:chi-pole-sum} is an identity, not a truncation. It is a
similarity transform of the same linear system, so it carries no
approximation beyond diagonalizability of $\mathcal L$.

For a defective eigenvalue, degeneracy alone is not the issue: a Jordan block
$J_\lambda=\lambda I+N$ of size $p$, with $N^p=0$, contributes
\begin{equation}
(J_\lambda+\ii\Omega I)^{-1}
=\sum_{r=0}^{p-1}\frac{(-1)^rN^r}{(\lambda+\ii\Omega)^{r+1}} .
\label{eq:jordan-resolvent}
\end{equation}
Thus the resolvent identity survives, but the simple Lorentzian pole sum does
not: defective blocks generate higher-order poles and should be evaluated by a
Jordan-free linear solve or contour/Riesz-projector construction in numerical
work.

\subsection{The trace-mode residue and the zero-frequency inverse}
\label{sec:zero-pole}

$\mathcal L$ always has the eigenvalue $\lambda_0=0$, whose right
eigenvector is the steady state $\rho_{\rm ss}$. Its term in
Eq.~\eqref{eq:resolvent-spectral} diverges as $\Omega\to0$. For a unique
stationary state and the traceless commutator drive used here, its trace-mode
residue vanishes; the zero-frequency inverse still requires an explicit
trace-zero or Drazin construction.

Trace preservation means $\Tr\mathcal L[X]=0$ for every $X$, which says
precisely that the trace functional is the \emph{left} eigenvector belonging
to $\lambda_0$:
\begin{equation}
\bra{l_0}=\Tr(\,\cdot\,).
\label{eq:left-null}
\end{equation}
Assume first that the pump Liouvillian has a unique stationary state. The drive
is a commutator, $\mathcal V_k\rho_{\rm ss}=-\ii[V_k,\rho_{\rm ss}]$, and is
therefore traceless. Hence
\begin{equation}
A^{(jk)}_0\propto\bra{l_0}\mathcal V_k\ket{\rho_{\rm ss}}
=\Tr\bigl(-\ii[V_k,\rho_{\rm ss}]\bigr)=0
\label{eq:zero-residue}
\end{equation}
identically, for any $V_k$ and any $\rho_{\rm ss}$. For $\Omega\ne0$ the
stationary-mode contribution to the contracted response therefore vanishes.
At $\Omega=0$, however, the full inverse does not exist and the formal
stationary term is $0/0$; one must retain the trace-zero solve of
Eq.~\eqref{eq:linear-resolvent}, omit the stationary subspace explicitly, or use
the Drazin inverse. If the Liouvillian has more than one stationary mode,
tracelessness alone does not eliminate the residues of every zero mode, and the
whole stationary subspace must be projected. Numerically, the trace-mode
residue vanishes exactly for the engine's Liouvillian and drives.

\subsection{Lorentzian reading}

Write the eigenvalues as
\begin{equation}
\lambda_m=-\Gamma_m+\ii\omega_m,\qquad \Gamma_m\geq0 ,
\label{eq:pole-parametrization}
\end{equation}
and absorb the sign into the residue, $B_m\equiv-A_m$. In the diagonalizable
simple-pole case, each term of
Eq.~\eqref{eq:chi-pole-sum} is then a complex Lorentzian in standard form,
\begin{equation}
\chibar_{jk}(\Omega)=\sum_m
\frac{B^{(jk)}_m}{\Gamma_m-\ii(\Omega+\omega_m)} ,
\label{eq:lorentzian-term}
\end{equation}
For an isolated pole with a slowly varying residue, the amplitude-squared
profile has HWHM $\Gamma_m$ and center
$\Omega=-\omega_m=-\Ima\lambda_m$. A complex residue mixes absorptive and
dispersive components, and overlapping poles need not display a peak exactly at
that center. Subject to those qualifications, the parametrization assigns a
direct physical reading to the numerically computed spectrum
of $\mathcal L$:
\begin{itemize}
\item $\Ima\lambda_m$ is a \emph{dressed oscillation frequency}: a
generalized Rabi frequency, an Autler--Townes or hyperfine splitting, or a
Raman detuning, depending on which coherence dominates $\ket{r_m}$.
\item $-\Rea\lambda_m$ is the corresponding \emph{power-broadened relaxation
rate}. It is not the bare $\Gamma$, $\gamma_{gg}$, or $\gamma_{\rm opt}$;
strong driving mixes the channels and shifts the widths.
\item The residues $A_m$ decide which poles are actually visible in a given
response element. A pole with small $\Gamma_m$ contributes a narrow feature
only if its residue is not also small.
\end{itemize}

\subsection{Positivity of the pole half-widths, and what it does not prove}
\label{sec:pole-positivity}

Because $\mathcal L$ is a Lindblad generator \cite{GKS1976}, the semigroup $\e^{\mathcal Lt}$ is
completely positive and trace preserving for all $t\geq0$. Its spectrum
therefore satisfies
\begin{equation}
\Rea\lambda_m\leq0\quad\text{for every }m,
\label{eq:gksl-spectrum}
\end{equation}
with at least one eigenvalue exactly zero, whose right eigenvector is the
steady state. Two consequences follow.

\emph{First}, Eq.~\eqref{eq:gksl-spectrum} places the poles in the closed stable
half plane. Finiteness on every real $\Omega$ additionally requires a unique
relaxing steady state, or an explicit check that every nonstationary mode has
$\Rea\lambda_m<0$; complete positivity alone permits peripheral imaginary
modes and decoherence-free sectors.

\emph{Second}, Eq.~\eqref{eq:gksl-spectrum} is guaranteed for an exactly
assembled Lindblad generator. A numerical spectral check is therefore an
implementation invariant rather than an independent validation of model
quality: it can detect assembly or truncation mistakes, but passing it does
not validate omitted physics. Direct evaluation confirms the invariant: at
both operating points of Table~\ref{tab:operating-points} the largest real
part found over the velocity grid is at most
$+1.7\times10^{-6}\,{\rm s^{-1}}$, which is the null eigenvalue resolved to
better than $5\times10^{-14}$ relative to the scale
$\Gamma=3.6\times10^{7}\,{\rm s^{-1}}$. Model errors --- including the
Floquet truncation error of Sec.~\ref{sec:floquet-convergence} --- must be
caught by convergence and comparison tests instead
(Sec.~\ref{sec:physicality-blind}).

\subsection{Effect of the velocity average}
\label{sec:pole-doppler}

Equation~\eqref{eq:chi-pole-sum} holds \emph{per velocity class}. The
velocity average of Sec.~\ref{sec:doppler} replaces the discrete pole set by
a distribution: each $\lambda_m(\bm v)$ traces a curve as $\bm v$ varies, and
the averaged response is an integral over that curve rather than a finite sum
of Lorentzians. This is the microscopic origin of inhomogeneous broadening;
only the simple-pole limit below reduces to a Voigt profile.

The two pole families behave very differently, because the Doppler shift
enters the excited-state diagonal as $\Delta_{\rm eff}=\Delta-k_{\rm pump}v_z$:
\begin{itemize}
\item \textbf{Optical poles} shift by the full one-photon rms scale
$k\sigma_v/2\pi=\sigma_v/\lambda$, about $244\MHz$ at $110\degC$ and
$247\MHz$ at $121\degC$, and are inhomogeneously smeared out.
\item \textbf{Ground-Raman poles} cancel that common axial shift in the
collinear limit. On the one-dimensional axial grid used below, their center
therefore stays near zero, while their \emph{width} varies through
velocity-dependent optical pumping; light shifts can produce residual center
shifts. The separate full non-collinear reference additionally distributes
their centers through the transverse term in
Eq.~\eqref{eq:doppler-detunings}; it does not produce the full optical smear.
\end{itemize}
Section~\ref{sec:pole-spectrum} quantifies both statements at the literature
operating point. The residual non-collinear distribution
$\sigma_{\delta,\theta}$ of Eq.~\eqref{eq:angular-doppler} convolves with the
velocity-dependent Raman response. The separate reference performs this
average, but a measured squeezing bandwidth additionally requires microscopic
noise propagation and detector transfer.

\section{Faddeeva base integral and Voigt identity}

For a single velocity projection,
\begin{equation}
\left\langle\frac1{x_0-kv+\ii\Gamma_o}\right\rangle_v
=-\frac{\ii}{k\sigma_v}\sqrt{\frac{\pi}{2}}\,
w\!\left(\frac{x_0+\ii\Gamma_o}{\sqrt2\,k\sigma_v}\right),
\qquad
w(z)=\e^{-z^2}\operatorname{erfc}(-\ii z).
\label{eq:faddeeva}
\end{equation}
With
\[
V(x;\sigma,\gamma)=
\frac{\Rea w[(x+\ii\gamma)/(\sqrt2\sigma)]}
{\sqrt{2\pi}\sigma},
\]
the dimensionally correct identity is
\begin{equation}
\boxed{\;
\Ima\left\langle\frac1{x_0-kv+\ii\Gamma_o}\right\rangle_v
=-\pi V(x_0;k\sigma_v,\Gamma_o).\;}
\label{eq:voigt-correct}
\end{equation}
There is no additional factor $1/k$. Equation~\eqref{eq:faddeeva} is an exact
base simple-pole integral, and it is precisely the velocity average of a
\emph{single} term of the pole sum \eqref{eq:chi-pole-sum} under the
assumption that the pole position is linear in $v$ and its residue is
constant. Velocity-dependent populations, Raman self-energies, and the
two-detuning geometry violate that assumption and generally require
multidimensional quadrature or a justified partial-fraction reduction. A
one-dimensional Faddeeva evaluation is not a complete treatment of the
experimental geometry.

%======================================================================
\part{Semiclassical field propagation and gain}
%======================================================================

\section{Classical coupled-mode equation}

Use the GABES positive-frequency convention
$E_j^{(+)}\propto\alpha_j(z)\e^{\ii k_{j,z}z-\ii\omega_jt}$ and define the
physical susceptibility so that $\Ima\chi_{jj}>0$ is passive absorption.
For $\Delta k_{\rm geom}=2k_{\rm pump}-k_p-k_c$, the symmetric mismatch
gauge is
$\widetilde\alpha_p=\alpha_p\exp(-\ii\Delta k_{\rm geom}z/2)$ and
$\widetilde\alpha_c^*=\alpha_c^*\exp(+\ii\Delta k_{\rm geom}z/2)$.
It gives diagonal phases $(-\ii\Delta k_{\rm geom}/2,
+\ii\Delta k_{\rm geom}/2)$ in the two-component basis
$(\widetilde\alpha_p,\widetilde\alpha_c^*)^T$. This fixes the code-level
Maxwell convention directly; in particular a passive uncoupled probe obeys
$|\alpha_p(L)|^2=\exp[-k_p\Ima\chi_{pp}L]|\alpha_p(0)|^2$.
Direct differentiation then gives
\begin{equation}
\frac{\dd}{\dd z}
\begin{pmatrix}\widetilde\alpha_p\\\widetilde\alpha_c^*\end{pmatrix}
=M_{\rm cl}
\begin{pmatrix}\widetilde\alpha_p\\\widetilde\alpha_c^*\end{pmatrix},
\qquad
M_{\rm cl}=Q^{-1}M_EQ,
\label{eq:classical-propagation}
\end{equation}
where the Maxwell electric-field matrix is
\begin{equation}
M_E=
\begin{pmatrix}
+\frac{\ii k_p}{2}\chi_{pp}-\frac{\ii\Delta k_{\rm geom}}2 &
+\frac{\ii k_p}{2}\chi_{pc}\\
-\frac{\ii k_c}{2}\chi_{cp}^* &
-\frac{\ii k_c}{2}\chi_{cc}^*+\frac{\ii\Delta k_{\rm geom}}2
\end{pmatrix},\qquad
Q=\operatorname{diag}\!\left(
\sqrt{\frac{2\hbar\omega_p}{\epsilon_0cA_p}},
\sqrt{\frac{2\hbar\omega_c}{\epsilon_0cA_c}}\right).
\label{eq:electric-propagation}
\end{equation}
Here
\begin{equation}
\Delta k_{\rm geom}=2k_{{\rm pump},z}-k_{p,z}-k_{c,z}
\label{eq:geometric-mismatch}
\end{equation}
uses vacuum carrier wave numbers and geometry only.
The effective mode areas $A_p,A_c$ must match the collected modes.
Equation~\eqref{eq:classical-propagation} is the canonical photon-flux
rescaling of the electric-field equation; a Rabi-amplitude implementation
must equivalently apply Eq.~\eqref{eq:canonical-rescaling}.

\subsection*{No double counting of dispersion}

Equation~\eqref{eq:classical-propagation} is Option A. The real parts of
$\chi_{pp}$ and $\chi_{cc}$ generate dispersive phase exactly once through the
diagonal propagation terms. Refractive indices
$n_j=1+\Rea\chi_{jj}/2$ are \emph{not} inserted again into
$\Delta k_{\rm geom}$. A dressed-$k$ calculation that uses both mechanisms is
retained only for implementation history and convergence diagnostics.

\section{Constant-matrix exponential and mean-field gains}
\label{sec:matrix-exponential}

For an undepleted pump and $z$-independent coefficients,
\begin{align}
T_{\rm cl}(L)&=\e^{M_{\rm cl}L}\nonumber\\
&=\e^{sL}\left[
\cosh(qL)I+\frac{\sinh(qL)}q(M_{\rm cl}-sI)\right],\\
s&=\frac{a+d}{2},\qquad
q=\sqrt{\frac{(a-d)^2}{4}+bc},
\label{eq:matrix-exponential}
\end{align}
where $a,b,c,d$ are the entries of $M_{\rm cl}$, so that $M_{\rm cl}$ has
eigenvalues
\begin{equation}
\mu_\pm=s\pm q .
\label{eq:Mcl-eigenvalues}
\end{equation}
Equation~\eqref{eq:matrix-exponential} needs no eigenvectors. Split off the
trace,
$M_{\rm cl}=sI+N$ with $N$ traceless. For a traceless $2\times2$ matrix the
Cayley--Hamilton theorem reads $N^2=-\det(N)\,I=q^2I$, so the exponential
series collapses into even and odd parts,
\begin{equation}
\e^{NL}=\sum_{k\ \rm even}\frac{(qL)^k}{k!}I
+\sum_{k\ \rm odd}\frac{(qL)^{k}}{k!}\frac{N}{q}
=\cosh(qL)\,I+\frac{\sinh(qL)}{q}N ,
\label{eq:cayley-hamilton}
\end{equation}
and $\e^{M_{\rm cl}L}=\e^{sL}\e^{NL}$ because $sI$ commutes with $N$. Since
$\sinh(z)/z$ is entire, Eq.~\eqref{eq:matrix-exponential} is valid verbatim
at $q=0$, where it reads $\e^{sL}(I+LN)$; no limit needs to be taken. The
same $2\times2$ trace--discriminant
parametrization reappears in the quantum-noise integral of
Sec.~\ref{sec:lyapunov}. Its numerical pair generally differs because
$M_q(\Omega)$ need not equal $M_{\rm cl}$.

The formula is an algebraic identity for constant $M_{\rm cl}$. Its physical
assumptions are:
\begin{itemize}
\item undepleted pump;
\item fixed density and temperature;
\item no longitudinal pump attenuation;
\item $z$-independent overlap and response coefficients; and
\item a single collected transverse mode.
\end{itemize}
If those assumptions fail, use an ordered product or path-ordered exponential.

\subsection{Gain convention}
\label{sec:gain-convention}

Both gains are referred to the incident probe-seed power $P_0$ and read
from a single photon-flux transfer matrix evaluated with probe-seed-only input:
\begin{equation}
\boxed{\;
G_p\equiv\frac{P_p^{\rm out}}{P_0}=|T_{{\rm cl},11}|^2,\qquad
G_c\equiv\frac{P_c^{\rm out}}{P_0}
=\frac{\omega_c}{\omega_p}|T_{{\rm cl},21}|^2 .\;}
\label{eq:mean-field-gains}
\end{equation}
This is a \textbf{two-gain convention}: $G_p$ and $G_c$ are computed
independently. Two gaps must be distinguished:
\begin{equation}
\mathcal G_{\rm flux}\equiv
|T_{{\rm cl},11}|^2-|T_{{\rm cl},21}|^2
=G_p-\frac{\omega_p}{\omega_c}G_c,
\qquad
\mathcal G_{\rm power}\equiv G_p-G_c .
\label{eq:gain-gap-observable}
\end{equation}
Both are outputs rather than inputs. Canonical commutator preservation constrains
$\mathcal G_{\rm flux}$, which equals one in the closed lossless parametric
limit of Eq.~\eqref{eq:gain-gap}. The power gap is numerically close because
$\omega_p\simeq\omega_c$, but it is not the canonical invariant.

The distinction matters because a near-degenerate single-gain power convention
is also in use. The exact photon-flux relation is
\begin{equation}
P_p=G\,P_0,\qquad
P_c=\frac{\omega_c}{\omega_p}(G-1)\,P_0 .
\label{eq:single-gain-convention}
\end{equation}
The commonly quoted $P_c=(G-1)P_0$ form additionally uses
$\omega_c/\omega_p\simeq1$. In either form the conjugate photon-flux gain has
been imposed rather than computed, and $\mathcal G_{\rm flux}=1$ by
construction. Equation~%
\eqref{eq:single-gain-convention} is the ideal-parametric special case of
Eq.~\eqref{eq:mean-field-gains}. The mapping between the two conventions is
\begin{equation}
G_p\;\longleftrightarrow\;G,\qquad
G_c\;\longleftrightarrow\;\frac{\omega_c}{\omega_p}(G-1) ,
\label{eq:gain-convention-map}
\end{equation}
valid only where $\mathcal G_{\rm flux}=1$. Two consequences follow. A
conjugate-to-probe-seed power ratio quoted under
Eq.~\eqref{eq:single-gain-convention} is not an independently calibrated
$G_c$, so it must not be compared with the computed $G_c$ as though it were.
The photon-flux gap also cannot serve as a validation target against a data set
in which it was assumed.

The quantities in Eq.~\eqref{eq:mean-field-gains} are explicitly
\textbf{semiclassical mean-field gains}. They do not by themselves determine
a quantum covariance.

\section{Classical photon-flux identity}

For $u'=au+bv$ and $v'=cu+dv$,
\begin{equation}
\frac{\dd}{\dd z}(|u|^2-|v|^2)
=2\Rea(a)|u|^2-2\Rea(d)|v|^2
+2\Rea[(b-c^*)u^*v].
\label{eq:gain-gap}
\end{equation}
In a closed lossless parametric model,
$\Rea a=\Rea d=0$ and $b=c^*$, so $\mathcal G_{\rm flux}=1$.
Equation~\eqref{eq:gain-gap} is a property of the classical transfer matrix.
If $\mathcal G_{\rm flux}\ne1$, it is evidence that a two-operator quantum
map needs additional channels; it is not permission to use a noncanonical
Bogoliubov map. Because the two-gain convention computes
$\mathcal G_{\rm flux}$ rather than imposing it, this test is available at
every scan point.

%======================================================================
\part{Quantum fluctuation propagation}
%======================================================================

\section{Quantum-Langevin propagation}

Define the two-mode sideband vector in the Bogoliubov (phase-conjugate)
pairing that the four-wave process couples,
\begin{equation}
\hat{\bm A}(z,\Omega)=
\begin{pmatrix}
\delta\hat a_p(z,\Omega)\\
\delta\hat a_c^\dagger(z,-\Omega)
\end{pmatrix},\qquad
[\hat{\bm A}(z,\Omega),\hat{\bm A}^\dagger(z,\Omega')]
=2\pi J\delta(\Omega-\Omega'),
\quad J=\operatorname{diag}(1,-1).
\label{eq:nambu}
\end{equation}
The metric $J$ is the one preserved by a two-mode Bogoliubov transformation;
a lossless parametric amplifier acts on $\hat{\bm A}$ by an element of
$SU(1,1)$, and the requirement $TJT^\dagger=J$ used below is exactly that
condition. Equivalently, in the quadrature basis of
Eq.~\eqref{eq:nambu-quadrature} it is the symplectic condition.
The distributed fluctuation equation is
\begin{equation}
\frac{\partial}{\partial z}\hat{\bm A}(z,\Omega_{\rm SA})
=M_q(z,\Omega_{\rm SA})\hat{\bm A}(z,\Omega_{\rm SA})
+B(z,\Omega_{\rm SA})\hat{\bm F}(z,\Omega_{\rm SA}).
\label{eq:langevin}
\end{equation}
Its formal solution is
\begin{align}
T(z_2,z_1;\Omega)&=
\mathcal P\exp\!\left[\int_{z_1}^{z_2}M_q(z,\Omega)\dd z\right],\\
\hat{\bm A}_{\rm out}(\Omega)&=
T(L,0;\Omega)\hat{\bm A}_{\rm in}(\Omega)
+\int_0^L T(L,z;\Omega)B(z,\Omega)\hat{\bm F}(z,\Omega)\dd z .
\label{eq:langevin-solution}
\end{align}
For local Markov reservoirs,
\[
[\hat{\bm F}(z,\Omega),\hat{\bm F}^\dagger(z',\Omega')]
=2\pi J_F\delta(z-z')\delta(\Omega-\Omega').
\]

\section{Preservation of the commutation relations}
\label{sec:realizability}

The output operators must obey the same commutation relations as the input
operators. With the reservoir contribution included this requires
\begin{equation}
\boxed{\;
TJT^\dagger+\int_0^L
T(L,z)BJ_FB^\dagger T^\dagger(L,z)\dd z=J ,\;}
\label{eq:commutator-integral}
\end{equation}
which is the statement that gain and loss must each be accompanied by their
own Langevin noise. A sufficient local condition is
\begin{equation}
\boxed{\;
M_qJ+JM_q^\dagger+BJ_FB^\dagger=0 .\;}
\label{eq:physical-realizability}
\end{equation}
Define the Hermitian commutator-defect matrix
\begin{equation}
K\equiv-\left(M_qJ+JM_q^\dagger\right)=BJ_FB^\dagger .
\label{eq:K-matrix}
\end{equation}
The signs of its eigenvalues classify the reservoir channels the drift
requires:
a positive eigenvalue of $K$ admits a loss-type (vacuum-absorbing) channel
with $J_F=+1$, while a negative eigenvalue requires a gain-type (inverted)
channel with $J_F=-1$. Both are physical for a parametric amplifier; the
sign pattern is a \emph{channel classification}, not a validity test. Because $K$
is defined through the canonical metric $J$, it is sensitive to the
normalization of Eq.~\eqref{eq:canonical-rescaling} and must be recomputed
whenever the collected mode areas or Rabi-to-field factors change.

\section{Covariance and atomic diffusion}

Define reservoir symmetrized correlations by
\[
\frac12\left\langle
\{\hat{\bm F}(z,\Omega),\hat{\bm F}^\dagger(z',\Omega')\}\right\rangle
=2\pi N_{\rm res}(z,\Omega)\delta(z-z')\delta(\Omega-\Omega'),
\qquad D=BN_{\rm res}B^\dagger .
\]
Then
\begin{equation}
\boxed{\;
V_{\rm out}(\Omega)=
T(\Omega)V_{\rm in}(\Omega)T^\dagger(\Omega)
+W(L;\Omega),\qquad
W(L;\Omega)\equiv\int_0^L T(L,z;\Omega)D(z,\Omega)T^\dagger(L,z;\Omega)\dd z .\;}
\label{eq:covariance}
\end{equation}
Equation~\eqref{eq:covariance} is one $2\times2$ two-mode block. Before applying
the detector model, solve its companion block and assemble
\[
\hat{\bm\Xi}(\Omega)=
\bigl(\hat a_p(\Omega),\hat a_c(\Omega),
\hat a_p^\dagger(-\Omega),\hat a_c^\dagger(-\Omega)\bigr)^T
\]
with spectral covariance $V_\Xi(\Omega)$. The quadrature-sideband covariance
used below is
\begin{equation}
V_R(\Omega)=CV_\Xi(\Omega)C^\dagger,\qquad
C=\frac1{\sqrt2}
\begin{pmatrix}
1&0&1&0\\
-\ii&0&\ii&0\\
0&1&0&1\\
0&-\ii&0&\ii
\end{pmatrix},
\label{eq:nambu-quadrature}
\end{equation}
corresponding to $\hat{\bm R}=(X_p,Y_p,X_c,Y_c)^T$. This explicit assembly
prevents the $2\times2$ two-mode covariance from being mixed dimensionally with
the $4\times4$ detector matrices.

The commutator condition fixes antisymmetric reservoir correlations, not the
symmetrized diffusion that determines squeezing. A physical
$D(z,\Omega_{\rm SA})$ must be derived from the pumped atomic Lindblad model,
generalized Einstein relations, or an independently justified reservoir
model. Independent velocity classes add as independent diffusion powers;
one must not square a velocity-averaged noise amplitude.

%----------------------------------------------------------------------
\section{Closed-form evaluation of the distributed-noise integral}
\label{sec:lyapunov}

The integral $W(L;\Omega)$ in Eq.~\eqref{eq:covariance} is the only part of
the fluctuation chain that is not already an algebraic operation. For a
$z$-independent drift it is available in closed form. Its $2\times2$
trace--discriminant structure parallels the mean-field transfer matrix, but it
is formed from $M_q(\Omega)$ itself.

\subsection{Differential Lyapunov form}

For constant $M_q$ the propagator depends only on the separation,
$T(L,z)=\e^{M_q(L-z)}$, so substituting $u=L-z$ gives
\begin{equation}
W(L)=\int_0^L \e^{M_qu}\,D\,\e^{M_q^\dagger u}\,\dd u .
\label{eq:W-integral}
\end{equation}
Introduce the \emph{Lyapunov superoperator}, a linear map on matrices,
\begin{equation}
\mathcal M[X]\equiv M_qX+XM_q^\dagger,
\qquad\text{so that}\qquad
\e^{u\mathcal M}[D]=\e^{M_qu}D\,\e^{M_q^\dagger u}.
\label{eq:lyapunov-superoperator}
\end{equation}
The second identity holds because left and right multiplication commute.
Equation~\eqref{eq:W-integral} is then simply $\int_0^L\e^{u\mathcal M}[D]\dd u$,
and differentiating with respect to the upper limit gives the differential
Lyapunov equation
\begin{equation}
\boxed{\;
\frac{\dd W}{\dd L}=\mathcal M[W]+D,\qquad W(0)=0 .\;}
\label{eq:differential-lyapunov}
\end{equation}
$W$ is Hermitian for Hermitian $D$, because $\mathcal M$ preserves
Hermiticity, and positive semidefinite for positive semidefinite $D$, because
the integrand of Eq.~\eqref{eq:W-integral} is a congruence of $D$. Both are
useful numerical invariants: a computed $W$ that violates either signals an implementation
error rather than a physical effect.

\subsection{Solution}
\label{sec:lyapunov-solution}

Equation~\eqref{eq:differential-lyapunov} is an inhomogeneous linear equation
with constant coefficients, so its solution is the elementary quadrature of
$\e^{u\mathcal M}$:
\begin{equation}
\boxed{\;
W(L)=L\,\varphi_1\!\left(L\mathcal M\right)[D],
\qquad
\varphi_1(z)\equiv\frac{\e^{z}-1}{z}
=\sum_{k\ge0}\frac{z^k}{(k+1)!}\; .}
\label{eq:lyapunov-closed-form}
\end{equation}
The coordinate-free form matters because $\varphi_1$ is an \emph{entire}
function --- the series on the right converges everywhere and
$\varphi_1(0)=1$ --- so Eq.~\eqref{eq:lyapunov-closed-form} needs no special
case, no removable-pole treatment, and no genericity assumption on $M_q$. The
mean-field transfer matrix of Sec.~\ref{sec:matrix-exponential} is built from
$\sinh(z)/z=\sum_{k\ge0}z^{2k}/(2k+1)!$, also entire and also equal to $1$ at
the origin, so neither layer has a singularity at a coalescing eigenvalue.
An eigenbasis can therefore fail even when the finite-$L$ matrix function is
regular. This numerical fact does not make an exceptional point unphysical:
eigenmode coalescence, non-normal amplification, and enhanced parameter
sensitivity can remain observable properties of the drift.

\subsection{Spectral resolution in trace--discriminant form}

To see \emph{which} exponent controls what, resolve
Eq.~\eqref{eq:lyapunov-closed-form} in the eigenbasis. If
$M_q=S\Lambda S^{-1}$ with $\Lambda=\operatorname{diag}(\mu_1,\mu_2)$, the
eigenmatrices of $\mathcal M$ are $S_{\cdot m}S_{\cdot n}^\dagger$ and
\begin{equation}
\operatorname{spec}\mathcal M=\{\mu_m+\mu_n^*\},
\label{eq:lyapunov-spectrum}
\end{equation}
so with $\widetilde D=S^{-1}DS^{-\dagger}$,
\begin{equation}
W(L)=S\,\widetilde W(L)\,S^\dagger,\qquad
\widetilde W_{mn}(L)=\widetilde D_{mn}\,
L\,\varphi_1\!\left[(\mu_m+\mu_n^*)L\right]
=\widetilde D_{mn}\,\frac{\e^{(\mu_m+\mu_n^*)L}-1}{\mu_m+\mu_n^*} .
\label{eq:lyapunov-eigenbasis}
\end{equation}
For the $2\times2$ two-mode block, define at each analysis frequency
\begin{equation}
s_q\equiv\frac{\operatorname{tr}M_q}{2},\qquad
q_q^2\equiv\frac{(a_q-d_q)^2}{4}+b_qc_q,\qquad
\mu_\pm=s_q\pm q_q ,
\label{eq:Mq-eigenvalues}
\end{equation}
where $a_q,b_q,c_q,d_q$ are the entries of $M_q(\Omega)$. Then
Eq.~\eqref{eq:lyapunov-spectrum} contains just three distinct values,
\begin{equation}
\mu_\pm+\mu_\pm^*=2\Rea(s_q\pm q_q),\qquad
\mu_\pm+\mu_\mp^*=2\Rea s_q\pm2\ii\,\Ima q_q .
\label{eq:lyapunov-denominators}
\end{equation}
Within the fluctuation problem, $\Rea(s_q\pm q_q)$ set how fast each
eigenchannel amplifies and accumulates diffusion, while $\Ima q_q$ sets the
interference between channels. These equal the classical $(s,q)$ only when
$M_q(\Omega)=M_{\rm cl}$. If $\Rea(\mu_m+\mu_n^*)<0$ for every pair,
$W(L)$ saturates as $L\to\infty$ at the algebraic Lyapunov solution
$\widetilde W_{mn}=-\widetilde D_{mn}/(\mu_m+\mu_n^*)$; in an amplifier at
least one exponent is positive and the noise grows with the gain, as it must.

\subsection{Numerically stable evaluation}
\label{sec:lyapunov-robust}

Equation~\eqref{eq:lyapunov-eigenbasis} is the form to \emph{reason} with,
but it is not the form to compute with. It requires $S$ to be invertible and
well conditioned, and the parametric drift $M_q$ is non-normal: as $q_q\to0$
its two eigenvectors coalesce and $S^{-1}$ diverges. The underlying $W(L)$
remains finite there, by Eq.~\eqref{eq:lyapunov-closed-form}; it is the
coordinate system that fails, not the solution. Direct testing bears this out:
the eigenbasis form loses all accuracy for a defective drift, while the
coordinate-free form agrees with direct quadrature to $2\times10^{-15}$ at the
same point (Table~\ref{tab:analytic-verification}).

Any stable evaluation of $\varphi_1$ therefore suffices. A convenient choice
is the block exponential \cite{VanLoan1978}: assemble
\begin{equation}
Z=\begin{pmatrix}-M_q&D\\0&M_q^\dagger\end{pmatrix},\qquad
\e^{ZL}=\begin{pmatrix}F_{11}&F_{12}\\0&F_{22}\end{pmatrix},
\label{eq:van-loan-block}
\end{equation}
from which
\begin{equation}
\boxed{\;W(L)=F_{22}^\dagger F_{12}\; .}
\label{eq:van-loan-result}
\end{equation}
Equation~\eqref{eq:van-loan-result} is Eq.~\eqref{eq:lyapunov-closed-form}
evaluated by a single $4\times4$ matrix exponential, with no eigenvector
inversion anywhere. It is the recommended implementation.

\subsection{Longitudinally varying drift}

When $M_q$ and $D$ vary with $z$ --- through pump depletion, density
gradients, or the Gaussian overlap profile --- no single closed form exists.
If $M_k$ and $D_k$ are the actual local coefficients, the segmented
composition is exact within each piecewise-constant segment. Writing
$T_k=\e^{M_k\dd z}$ and $W_k=W(\dd z;M_k,D_k)$ for segment $k$, the
covariance accumulates as
\begin{equation}
V^{(k+1)}=T_kV^{(k)}T_k^\dagger+W_k ,
\label{eq:segmented-covariance}
\end{equation}
with each $W_k$ evaluated by Eq.~\eqref{eq:van-loan-result}. This is the
noise-carrying counterpart of the segmented mean-field propagation and
introduces no approximation beyond piecewise-constant coefficients.

\subsection{Static independent-vacuum completion fixture}

For testing only, any constant drift may be given this particular local
commutator completion. Using $K$ from Eq.~\eqref{eq:K-matrix}, let
\begin{equation}
K=U\Lambda_KU^\dagger,\quad
B_{\rm vac}=U|\Lambda_K|^{1/2},\quad
J_F=\operatorname{diag}(\operatorname{sgn}\Lambda_K),\quad
D_{\rm vac}=\tfrac12B_{\rm vac}B_{\rm vac}^\dagger .
\label{eq:minimal-dilation}
\end{equation}
This eigenfactor closes each channel of $K$ with one independent vacuum port
and preserves the commutator by construction. The subscript ``vac'' denotes
this chosen independent-vacuum algebraic fixture; it is not a universal lower
bound over microscopic atomic reservoirs. In the special phase-insensitive amplifier
setting, an analogous independent-vacuum construction gives the familiar
Caves added-noise bound \cite{Caves1982}, but an arbitrary dissipative
two-mode drift need not satisfy those assumptions. The fixture is not a
derivation of atomic spontaneous-emission noise, and its symmetrized diffusion
cannot be promoted to a physical bound without a microscopic model or data. When
$M_q$ and the diagnostic detector model are frequency independent,
$D_{\rm vac}$ is also frequency independent, and feeding it through
Eqs.~\eqref{eq:lyapunov-closed-form} and \eqref{eq:spectrum} yields a flat
$S_-$; the machinery is exercised, but no bandwidth is predicted.

\section{Present limitation of the frequency-dependent observable}

The final in-cell observable must be a spectrum
$S_-^{\rm out}(\Omega_{\rm SA})$, not a scalar. The algebraic path from
$M_q(\Omega_{\rm SA})$ and $D(\Omega_{\rm SA})$ to that spectrum is now
complete, but the microscopic $M_q(\Omega_{\rm SA})$ and
$D(\Omega_{\rm SA})$ are not implemented. The static vacuum completion is
paired here with a static detector model and yields only a flat diagnostic
curve. It cannot reproduce:
\begin{itemize}
\item the squeezing bandwidth;
\item low-frequency excess noise;
\item the noise peak near $2\MHz$ from laser-source electronic noise;
\item frequency-dependent atomic Langevin noise;
\item slow-light and group-delay effects; or
\item the frequency response of the detector.
\end{itemize}

%======================================================================
\part{External detection model}
%======================================================================

\section{External loss only after the cell}

In the quadrature basis
$\hat{\bm R}=(X_p,Y_p,X_c,Y_c)^T$, define
\[
K_\eta=\operatorname{diag}(\sqrt{\eta_p},\sqrt{\eta_p},
\sqrt{\eta_c},\sqrt{\eta_c}),\qquad
K_\ell=\operatorname{diag}(\sqrt{1-\eta_p},\sqrt{1-\eta_p},
\sqrt{1-\eta_c},\sqrt{1-\eta_c}).
\]
Only loss physically located after the cell is applied:
\begin{equation}
V_{R,\rm det}=K_\eta V_{R,\rm out}K_\eta^T
+K_\ell V_{\rm vac}K_\ell^T .
\label{eq:detection-loss}
\end{equation}
For symmetric efficiency,
\[
V_{R,\rm det}=\eta_{\rm ext}V_{R,\rm out}
+(1-\eta_{\rm ext})V_{\rm vac}.
\]
Atomic absorption and spontaneous emission remain distributed inside
Eqs.~\eqref{eq:langevin}--\eqref{eq:covariance}. Cell optical-depth factors
$\exp(-{\rm OD}_{p,c,\rm seg})$ are not applied a second time after
propagation.

\section{Electronic weighting and SQL normalization}

The linearized photocurrent is
\begin{equation}
\delta\hat i_j(\Omega)\propto H_j(\Omega)
\left[\alpha_j^*\delta\hat a_j(\Omega)
+\alpha_j\delta\hat a_j^\dagger(-\Omega)\right],
\qquad
\delta\hat i_-=\delta\hat i_p-w(\Omega)\delta\hat i_c .
\label{eq:photocurrent}
\end{equation}
The normalized spectrum is
\begin{equation}
\boxed{\;
S_-(\Omega)=
\frac{\bm m^\dagger(\Omega)V_{R,\rm det}(\Omega)\bm m(\Omega)
+S_{\rm el}(\Omega)}
{S_{\rm SQL}^{(\rm exp)}(\Omega)} .\;}
\label{eq:spectrum}
\end{equation}
Here $\bm m(\Omega)$ is the four-component measurement vector obtained from
Eq.~\eqref{eq:photocurrent}; it contains the detected carrier amplitudes,
$H_p,H_c$, their phases, and the factor $-w$ on the conjugate quadratures.
Equations~\eqref{eq:covariance}, \eqref{eq:nambu-quadrature}, and
\eqref{eq:van-loan-result} supply the in-cell block
$V_{\rm out}=T V_{\rm in}T^\dagger+W$ and hence
$V_{R,\rm out}=CV_\Xi C^\dagger$. Thus Eq.~\eqref{eq:spectrum} is the single
assembled observable: once the microscopic
$M_q(\Omega_{\rm SA})$, $D(\Omega_{\rm SA})$, and measured detector functions
are supplied, every remaining step is algebraic.
With $V_{\rm vac}=I/2$ and independent detector vacua, the corresponding
ideal electronic standard-quantum-limit (SQL) term is proportional to
\[
S_{\rm SQL}(\Omega)\propto
|H_p|^2\eta_p|\alpha_p|^2+
w^2|H_c|^2\eta_c|\alpha_c|^2 ,
\]
before the experiment-specific electronic calibration.
The SQL reference must use the same detected powers, electronic
transfer functions, weight $w$, resolution bandwidth, and video bandwidth.
The report must state whether electronic dark noise is included or subtracted.

\begin{table}[ht]
\centering
\caption{Loss and noise ledger.}
\label{tab:loss-ledger}
\small
\begin{tabularx}{\textwidth}{L{3.4cm}L{3.2cm}Y}
\toprule
Mechanism & Physical location & Correct treatment\\
\midrule
Atomic absorption and spontaneous emission & distributed in the vapor &
Complex response in $M_q$ plus associated $B$ and $D$\\
Cell windows, filters, and free propagation & after the cell &
External beam-splitter efficiency $\eta_{\rm path}$\\
Detector quantum efficiency & detector face &
External beam-splitter efficiency $\eta_{\rm QE}$\\
Electronic imbalance and bandwidth & after photodetection &
$H_p(\Omega)$, $H_c(\Omega)$, $w(\Omega)$, and calibrated
$S_{\rm el}(\Omega)$\\
Pump scattering into the detected mode & apparatus dependent &
Measured technical-noise spectrum; not a universal atomic-loss formula\\
\bottomrule
\end{tabularx}
\end{table}

The ansatz
\[
\mathcal N_{\rm ps}=\kappa(1-\e^{-{\rm OD}_{\rm pump}})
\]
is phenomenological. If used, it must be written
$\mathcal N_{\rm ps}(\Omega)=
\kappa(1-\e^{-{\rm OD}_{\rm pump}})F_{\rm ps}(\Omega)$ and calibrated to the
experimental spectral normalization. It is not used as a universal Langevin
noise law in this report.

%======================================================================
\part{Validation}
%======================================================================

\section{Validation hierarchy}

Agreement between this algebra and the production program is an algebraic and
numerical consistency check, not independent physical validation. Independence
is local rather than global: the continued-fraction code independently assembles
the archived finite-Floquet blocks, while the literal SI reference independently
reconstructs Maxwell/$Q$ propagation from a shared reduced susceptibility. It
does not independently reproduce the atomic Hamiltonian or dissipator. The
separate pump reference checks its trace-constrained null vector against an SVD
construction, its weak response against an independent dense/direct solve, and
its rotating frame against the finite-Floquet state, while still sharing the
same four-level atomic assembly. The levels below are therefore reported
separately.

\begin{table}[ht]
\centering
\caption{Validation hierarchy. No combined error column is used.}
\label{tab:validation-hierarchy}
\small
\begin{tabularx}{\textwidth}{L{3.4cm}Y L{3.0cm}}
\toprule
Level & Test performed & Status\\
\midrule
Algebraic consistency & Closed-form matrix exponential versus direct
exponential; dimensional and frequency audits & \pass\\
Analytic-identity verification & Spectral resolvent expansion, Lyapunov
closed form, block-exponential form, and the cubic characteristic polynomial
against direct numerics & \pass\ (Table~\ref{tab:analytic-verification})\\
Numerical implementation consistency & Continued-fraction versus direct
finite-block assembly for the archived reference, plus production Maxwell/$Q$
versus a literal SI propagation reference & \pass\ within shared atomic
assumptions\\
Pump-only weak-response reference & Minus-branch static-frame/Floquet gauge
parity, trace-one state, trace-zero finite-frequency and projected-DC response,
finite-seed limit, and direct/pole-residue parity & \pass\ numerically within
the four-level, one-dimensional-velocity model; separate slow reference,
plus branch unsupported\\
Two-dimensional Raman-Doppler reference & Independent $(v_z,v_x)$ Maxwell
quadrature, lab-frequency invariance, analytic rms, collinear/equal-$k$ limits,
grid refinement, and cutoff refinement & \pass\ for the separate slow
minus-branch reference; production remains one-dimensional\\
Truncation convergence & $N_F=1,2,3$ direct/continued/kernel parity and
scan-level adjacent-order checks & Production $N_F=3/2$ gate implemented;
historical $N_F=1$ \fail\\
Independent microscopic calculation & Multilevel Zeeman-resolved
Liouvillian and atomic diffusion matrix & \pending\\
Experimental gain comparison & Reduced model evaluated at the documented
$^{85}$Rb operating point & performed; substantial discrepancy\\
Experimental bandwidth comparison & Full $S_-(\Omega_{\rm SA})$ versus the
measured $3.5\MHz$ squeezing bandwidth & two-dimensional mean-field reference
complete only for the declared minus branch, angle, and quadrature; microscopic
diffusion and detector response pending\\
Experimental squeezing-spectrum comparison & Full
$S_-(\Omega_{\rm SA})$ using microscopic $D(\Omega_{\rm SA})$ &
\pending\\
\bottomrule
\end{tabularx}
\end{table}

\section{Verification of the analytic identities}
\label{sec:analytic-verification}

The core closed forms used in the response, propagation, and noise layers were
checked against independent direct computations before use.
Table~\ref{tab:analytic-verification} records the residuals. The resolvent test
uses the engine's own $16\times16$ Liouvillian
at the literature operating point; the Lyapunov tests use randomized
non-normal drifts together with deliberately constructed defective cases.

\begin{table}[ht]
\centering
\caption{Verification of the analytic results against direct numerics.
Relative errors are Frobenius norms unless stated otherwise.}
\label{tab:analytic-verification}
\small
\begin{tabularx}{\textwidth}{L{4.6cm}L{4.4cm}C{2.4cm}Y}
\toprule
Identity & Reference computation & Relative error & Verdict\\
\midrule
Spectral resolvent expansion,
Eq.~\eqref{eq:resolvent-spectral} &
Direct linear solve on the $16\times16$ engine Liouvillian,
$\Omega_{\rm SA}/2\pi=0.1$--$4\MHz$ &
$1.0\times10^{-12}$ & \pass\\
Vanishing of the steady-state residue,
Eq.~\eqref{eq:zero-residue} &
$\bra{l_0}\mathcal L$ and
$\bra{l_0}\mathcal V_k\ket{\rho_{\rm ss}}$ on the engine Liouvillian &
exactly zero in floating-point arithmetic & \pass\\
Coordinate-free solution,
Eq.~\eqref{eq:lyapunov-closed-form} &
Simpson quadrature of Eq.~\eqref{eq:W-integral}; generic, defective, and
coalescing drifts & $2.2\times10^{-15}$ (worst) & \pass\\
Eigenbasis resolution,
Eq.~\eqref{eq:lyapunov-eigenbasis} &
Same quadrature, randomized non-normal $M_q$, PSD $D$ &
$5.7\times10^{-15}$ & \pass\\
Hermiticity and positivity of $W$ &
Same, invariants of Eq.~\eqref{eq:differential-lyapunov} &
$1.0\times10^{-16}$; $\min\operatorname{eig}W>0$ & \pass\\
Eigenbasis resolution at a defective drift &
Same quadrature, exact Jordan block &
loses all accuracy & \fail\ (basis)\\
Eigenbasis resolution near an exceptional point &
Same quadrature, coalescing eigenvalues &
$1.7\times10^{-1}$ & \fail\ (basis)\\
Lyapunov superoperator spectrum,
Eq.~\eqref{eq:lyapunov-spectrum} &
$\operatorname{spec}\mathcal M$ versus $\{\mu_m+\mu_n^*\}$, 200 random cases
up to $4\times4$ & $1.5\times10^{-13}$ & \pass\\
Block-exponential evaluation,
Eq.~\eqref{eq:van-loan-result} &
Same quadrature, generic and defective drifts &
$8.5\times10^{-15}$ (worst) & \pass\\
Cayley--Hamilton collapse,
Eq.~\eqref{eq:cayley-hamilton} &
$\|N^2-q^2I\|$ and the resulting $\e^{M_{\rm cl}L}$ versus a direct
exponential, including $q=0$ & $1.5\times10^{-14}$ & \pass\\
Characteristic polynomial,
Eq.~\eqref{eq:torrey-cubic} &
Direct eigenvalues of Eq.~\eqref{eq:bloch-generator} &
$1.7\times10^{-15}$ & \pass\\
Undamped limit,
Eq.~\eqref{eq:generalized-rabi} &
Roots versus $\Omega'=\sqrt{\Omega^2+\Delta^2}$ &
exact to machine precision & \pass\\
\bottomrule
\end{tabularx}
\end{table}

The two \fail\ rows are labeled ``basis'' because they are not failures of
the solution. The same quantity computed from
Eq.~\eqref{eq:lyapunov-closed-form}, at the same drift, is accurate to
$2\times10^{-15}$; only the eigenbasis resolution
\eqref{eq:lyapunov-eigenbasis} degrades, and it degrades because $S^{-1}$
does. That is the practical content of the remark in
Sec.~\ref{sec:lyapunov-solution} that $\varphi_1$ is entire: there is nothing
to regularize, only a coordinate system to avoid. A production
implementation that diagonalizes $M_q$ without a conditioning guard will
silently return invalid results in exactly the near-degenerate regime where
the two parametric eigenchannels are most strongly coupled.

\section{Atomic pole spectrum at the operating point}
\label{sec:pole-spectrum}

Applying Eqs.~\eqref{eq:resolvent-spectral}--\eqref{eq:pole-parametrization}
to the literature operating point
($\Delta/2\pi=+0.90\GHz$, $\delta/2\pi=-8\MHz$, $T=121\degC$,
$\Omega_{\rm pump}/2\pi=0.7076\GHz$) gives the pole structure in
Table~\ref{tab:pole-spectrum}.

\begin{table}[ht]
\centering
\caption{Slowest Liouvillian poles at the literature operating point,
zero-velocity class. Half-width is $-\Rea\lambda_m/2\pi$ and center is
$-\Ima\lambda_m/2\pi$. The Doppler column is the range of the half-width
across the one-dimensional axial velocity grid.}
\label{tab:pole-spectrum}
\small
\begin{tabularx}{\textwidth}{C{2.3cm}C{2.6cm}C{3.4cm}Y}
\toprule
Half-width [MHz] & Center [MHz] &
Doppler range of half-width [MHz] & Assignment\\
\midrule
$0$ (exact) & $0$ & --- &
Steady state; null eigenvalue of Eq.~\eqref{eq:pump-nullspace}\\
$0.5201$ & $0$ & $0.220$--$1.494$ &
Ground-state Raman coherence, power broadened\\
$0.5894$ & $\pm190.05$ & $0.311$--$1.494$ &
Excited-state hyperfine coherence\\
$3.3024$ & $\pm1023.57$ & $0.311$--$1.549$ &
Optical/dressed coherence\\
\bottomrule
\end{tabularx}
\end{table}

\paragraph{Power broadening widens the Raman pole fivefold.}
The slowest nonvanishing pole has half-width $0.5201\MHz$, whereas the
assumed ground-coherence rate is $\gamma_{gg}/2\pi=0.1029\MHz$ from
Eq.~\eqref{eq:rates-at-literature}. The pole is broadened by a factor
$5.06$. Any estimate that identifies the Raman feature width with
$\gamma_{gg}$ is therefore low by a factor of $5.06$ at this pump power. This
is the quantitative content of the statement in
Sec.~\ref{sec:pole-expansion} that $-\Rea\lambda_m$ is not a bare rate.

\paragraph{A narrow Raman feature remains in the one-dimensional axial average.}
The averaged response is not itself a finite pole sum. Rather, on the
one-dimensional axial grid used for Table~\ref{tab:pole-spectrum}, the
per-velocity Raman-pole centers remain at zero while their half-widths vary over
$0.220$--$1.494\MHz$: the two-photon combination cancels the common axial
Doppler shift. This does \emph{not} include the transverse term of
Eq.~\eqref{eq:doppler-detunings}, which supplies the $1.38\MHz$ rms center
spread in the full two-dimensional geometry. The separate reference in
Table~\ref{tab:angular-doppler-convergence} includes that term; the pole table
itself remains one-dimensional. The Raman feature is therefore broadened but
not smeared over the $247\MHz$ optical scale.

\paragraph{The separate scales do not determine a squeezing bandwidth.}
The per-velocity Raman poles supply Lorentzian half-widths of
$0.22$--$1.49\MHz$, while the non-collinear term supplies a Gaussian rms center
spread of $1.38\MHz$. These quantities have different definitions: a
Lorentzian has no finite variance, and the maximum half-width of one velocity
class is not the half-width of the averaged line. They therefore cannot be
combined by a root-sum-of-squares rule. The measured
$3.5\MHz$ bandwidth belongs to $S_-(\Omega_{\rm SA})$ after propagation and
detection, not to a bare atomic pole. The two-dimensional mean-field response
is now available as a separate reference, but a quantitative comparison still
requires frequency-dependent $M_q(\Omega_{\rm SA})$ and
$D(\Omega_{\rm SA})$ plus detector transfer functions. Until those are
supplied, the pole table is a scale diagnosis and no numerical bandwidth gap
is claimed.

\section{Floquet-truncation convergence}
\label{sec:floquet-convergence}

Table~\ref{tab:floquet-convergence} repeats the archived-model calculation at
the legacy numerical optimum
$\Delta/(2\pi)=-1.50\GHz$, $T=110\degC$ and approximately
$\delta/(2\pi)=-280\MHz$. It deliberately retains the archived
dressed-wave-vector phase convention only to isolate Floquet convergence;
because that convention double counts dispersion, these numbers are not the
Option-A physical prediction. To prevent their legacy columns from being
mistaken for canonical gains, the next two tables use
$g_{c,{\rm hist}}\equiv|T_{21}|^2$ and
$\Delta_{\rm hist}\equiv G_p-g_{c,{\rm hist}}$. Only
their relative truncation behavior is used.

\begin{table}[ht]
\centering
\caption{Floquet convergence of the archived reduced model at the fixed point
$\delta/(2\pi)=-280\MHz$. Sectors $n=-N_F,\ldots,N_F$ are retained.}
\label{tab:floquet-convergence}
\small
\begin{tabular}{c r r r r r}
\toprule
$N_F$ & $G_p$ & $g_{c,{\rm hist}}$ & $\Delta_{\rm hist}$ &
$\arg\chi_{pc}$ & $\arg\chi_{cp}$\\
\midrule
1 & 22.5913 & 21.9687 & $+0.622609$ &
$178.142^\circ$ & $-179.333^\circ$\\
2 & 19.3415 & 19.4472 & $-0.105684$ &
$178.173^\circ$ & $-179.366^\circ$\\
3 & 19.3415 & 19.4472 & $-0.105684$ &
$178.173^\circ$ & $-179.366^\circ$\\
\bottomrule
\end{tabular}
\end{table}

\begin{table}[ht]
\centering
\caption{Separate minimization of the gain-only objective. All
observables in a row are evaluated at that row's $\delta^*$. The scan spacing
is $5\MHz$.}
\label{tab:legacy-delta-star}
\small
\begin{tabular}{c r r r r r r}
\toprule
$N_F$ & $\delta^*/(2\pi)$ [MHz] & $G_p$ & $g_{c,{\rm hist}}$ &
$\Delta_{\rm hist}$ &
$\arg\chi_{pc}$ & $\arg\chi_{cp}$\\
\midrule
1 & $-280$ & 22.5913 & 21.9687 & $0.622609$ &
$178.142^\circ$ & $-179.333^\circ$\\
2 & $-275$ & 8.97871 & 8.46840 & $0.510314$ &
$178.402^\circ$ & $-179.582^\circ$\\
3 & $-275$ & 8.97871 & 8.46840 & $0.510314$ &
$178.402^\circ$ & $-179.582^\circ$\\
\bottomrule
\end{tabular}
\end{table}

The relative changes from $N_F=1$ to 2 are
\[
\frac{|G_p^{(2)}-G_p^{(1)}|}{G_p^{(2)}}=16.80\%,
\qquad
\frac{|g_{c,{\rm hist}}^{(2)}-g_{c,{\rm hist}}^{(1)}|}
{g_{c,{\rm hist}}^{(2)}}=12.97\%.
\]
The historical coefficient gap even changes sign. Thus the $N_F=1$ result is
\emph{qualitative and unconverged}. The changes from $N_F=2$ to $N_F=3$ are
below $4\times10^{-7}\%$ at the tested point. Direct block solves and the
independent assembly agree to $5.4\times10^{-13}$ or better. Convergence at
one operating point does not prove global convergence; this test must be
repeated for every reported scan. The production path now performs that
full-scan $N_F/N_F-1$ comparison and returns \texttt{UNCONVERGED} rather than
silently promoting a failed scan.

For the propagation-only Option-A correction at the fixed literature point,
the inherited approximate atomic response gives the results in
Table~\ref{tab:floquet-option-a}. Here $N_F=1\to2$ changes $G_p$ by
$0.4528\%$ and $G_c$ by $0.4906\%$; $N_F=2$ and 3 agree at the displayed
precision. This establishes Floquet convergence of that fixed finite-seed
propagation benchmark. It does not use the pump-only reference; the comparison
with that reference is reported separately below. The multidimensional-velocity
approximation remains untouched. The displayed gains assume equal collected mode
areas, but the conversion is explicit:
$T_{\rm can}=Q^{-1}T_EQ$, $G_c=(\omega_c/\omega_p)|T_{{\rm can},21}|^2$, and
$\mathcal G_{\rm flux}=G_p-|T_{{\rm can},21}|^2$.
\begin{table}[ht]
\centering
\caption{Floquet convergence at the fixed literature point
$\Delta/(2\pi)=+0.90\GHz$, $\delta/(2\pi)=-8\MHz$, and $T=121\degC$.
This is not an optimization; it is the production-like finite-seed path, not
the separate pump-only reference.}
\label{tab:floquet-option-a}
\small
\begin{tabular}{c r r r r r r}
\toprule
$N_F$ & $G_p$ & $G_c$ & $g_c^{(\Phi)}$ & $\mathcal G_{\rm flux}$ &
$\arg\chi_{pc}$ & $\arg\chi_{cp}$\\
\midrule
1 & 401.70833 & 403.86243 & 403.85591 & $-2.147574$ &
$-179.886152^\circ$ & $179.814904^\circ$\\
2 & 403.53550 & 405.85359 & 405.84704 & $-2.311543$ &
$-179.885623^\circ$ & $179.814496^\circ$\\
3 & 403.53550 & 405.85359 & 405.84704 & $-2.311543$ &
$-179.885623^\circ$ & $179.814496^\circ$\\
\bottomrule
\end{tabular}
\end{table}

The inherited weak-response implementation also retains a finite reference
probe-seed. Its numerical linearity check is
\begin{center}
\begin{tabular}{c r r r}
\toprule
Reference probe-seed power & $G_p$ & $G_c$ & $\mathcal G_{\rm flux}$\\
\midrule
$2\,\mu{\rm W}$ & 403.752756 & 406.102206 & $-2.342895$\\
$4\,\mu{\rm W}$ & 403.680319 & 406.019313 & $-2.332440$\\
$8\,\mu{\rm W}$ & 403.535497 & 405.853591 & $-2.311543$\\
\bottomrule
\end{tabular}
\end{center}
Between 2 and $8\,\mu{\rm W}$, $G_p$ changes by $-0.05381\%$, $G_c$ by
$-0.06122\%$, and the signed flux gap changes by $-1.3382\%$ relative to its
$2\,\mu{\rm W}$ value. The absolute gains are weakly seed dependent, while
the noncanonical flux gap is relatively sensitive. This supports numerical
linearity with respect to the reference probe-seed at this fixed point only;
by itself it does not validate the underlying pump-state approximation.

\subsection{Pump-only reference and the finite-seed limit}
\label{sec:pump-reference-validation}

The independent reference closes that specific numerical gap for the standard
minus Raman branch without replacing the production solver. On the complete
configured one-dimensional Doppler grid, four analysis frequencies
$\Omega_{\rm SA}/2\pi=(0,0.1,1,4)\MHz$ and five two-photon detunings
$\delta/2\pi=(-30,-8,0,12,40)\MHz$ give the diagnostics in
Table~\ref{tab:pump-reference-diagnostics}.

\begin{table}[ht]
\centering
\caption{Executable pump-only state and trace-zero weak-response diagnostics.}
\label{tab:pump-reference-diagnostics}
\small
\begin{tabularx}{\textwidth}{Y r}
\toprule
Check & Maximum error or margin\\
\midrule
Normalized pump-state equation residual & $4.00\times10^{-17}$\\
Pump-state trace error & $2.89\times10^{-15}$\\
Minimum pump-state eigenvalue & $2.371\times10^{-3}$\\
Normalized weak-response equation residual, full $4\times5$ scan &
$7.32\times10^{-17}$\\
Weak-response trace error, full $4\times5$ scan & $6.48\times10^{-25}$\\
Static-frame versus $N_F=3$ pump-only Floquet state &
$8.37\times10^{-15}$\\
Projected zero-relative-frequency response residual & $1.79\times10^{-15}$\\
Direct resolvent versus pole/residue sum, $0$--$4\MHz$ analysis band &
$2.45\times10^{-24}\,\mathrm{s}$\\
Pole eigenvector condition number & $3.39$\\
\bottomrule
\end{tabularx}
\end{table}

The finite-reference susceptibility was then compared with the infinitesimal
derivative after Doppler averaging all four complex response components over
the same five detunings. Table~\ref{tab:finite-seed-reference-limit} reports the
worst component error normalized by the maximum pump-reference response.

\begin{table}[ht]
\centering
\caption{Finite-seed $N_F=3$ approach to the pump-only infinitesimal reference.}
\label{tab:finite-seed-reference-limit}
\small
\begin{tabular}{c r r}
\toprule
Rabi-amplitude fraction & Seed power ($\mu\mathrm W$) & Worst normalized error\\
\midrule
$1.00$ & $8.0000$ & $2.503\times10^{-3}$\\
$0.30$ & $0.7200$ & $5.093\times10^{-4}$\\
$0.10$ & $0.0800$ & $9.254\times10^{-5}$\\
$0.03$ & $0.0072$ & $9.138\times10^{-6}$\\
\bottomrule
\end{tabular}
\end{table}

The errors decrease monotonically, and a log--log fit to the last three rows
has Rabi-amplitude slope $1.749$. The finite range therefore supports
convergence to the infinitesimal reference but does \emph{not} establish an
exact quadratic law. At the smallest amplitude, $N_F=1$ gives a
$4.485\times10^{-3}$ worst error, whereas $N_F=2$ and 3 both give
$9.138\times10^{-6}$; seed linearization and Floquet truncation must therefore
be checked separately. The direct resolvent and diagonalizable pole/residue
sum agree to $2.45\times10^{-24}\,\mathrm{s}$ or better over the declared
analysis band, with eigenvector condition number $3.39$. The projected
zero-relative-frequency fixture has normalized residual
$1.79\times10^{-15}$.

\section{Physicality tests do not detect the Floquet truncation error}
\label{sec:physicality-blind}

Since $\rho_0$ is the time average of a physical $\rho(t)$, it must be
Hermitian and positive semidefinite, with unit trace. Table~\ref{tab:physicality}
checks those necessary conditions for the $N_F=1$ truncation at both operating
points, after the convergence failure above has already been quantified.

\begin{table}[ht]
\centering
\caption{Physicality of the $N_F=1$ truncated Floquet state. Minima are taken
over the velocity grid and, for $\rho(t)$, over one beat period.}
\label{tab:physicality}
\small
\begin{tabularx}{\textwidth}{L{4.4cm}C{3.0cm}C{3.0cm}Y}
\toprule
Quantity & Archived point & Literature point & Requirement\\
\midrule
$\max\|\rho_0-\rho_0^\dagger\|_\infty$ &
$7.4\times10^{-14}$ & $3.3\times10^{-14}$ & $0$\\
$\max|\Tr\rho_0-1|$ &
$1.1\times10^{-15}$ & $5.6\times10^{-15}$ & $0$\\
$\min\operatorname{eig}\rho_0$ &
$+7.90\times10^{-3}$ & $+2.42\times10^{-3}$ & $\geq0$\\
$\min_t\operatorname{eig}\rho(t)$ &
$+5.63\times10^{-3}$ & $+2.28\times10^{-3}$ & $\geq0$\\
\bottomrule
\end{tabularx}
\end{table}

Every test passes, including at the archived point where
Table~\ref{tab:floquet-convergence} shows a $16.8\%$ error in $G_p$ and a sign
reversal of the historical transfer-coefficient gap. State positivity is
therefore insensitive to this truncation error: a valid density matrix can be
quantitatively wrong. The exact-GKSL half-plane condition of
Sec.~\ref{sec:pole-positivity} is an assembly invariant for the same reason,
so neither check substitutes for explicit $N_F$ convergence.

\clearpage
\section{Velocity-integration convergence}

The following tests retain the same historical propagation columns and use
$N_F=3$ with the one-dimensional Doppler integral. They validate only that
quadrature, not the propagation convention or non-collinear geometry.
At a fixed $5\sigma_v$ cutoff, velocity-step convergence is excellent:

\begin{table}[ht]
\centering
\caption{One-dimensional velocity-step convergence at fixed
$|v_z|\leq5\sigma_v$.}
\label{tab:velocity-step}
\small
\begin{tabular}{r r r r r}
\toprule
$\Delta v$ [m/s] & Classes & $G_p$ & $g_{c,{\rm hist}}$ &
Relative $G_p$ error\\
\midrule
10.00 & 195 & 18.936590 & 19.024328 & $1.56\times10^{-3}\%$\\
5.00  & 389 & 18.936295 & 19.024021 & $2.66\times10^{-7}\%$\\
2.50  & 777 & 18.936295 & 19.024021 & $9.58\times10^{-8}\%$\\
1.25  & 1551 & 18.936295 & 19.024021 & reference\\
\bottomrule
\end{tabular}
\end{table}

The cutoff, however, is important:

\begin{table}[ht]
\centering
\caption{One-dimensional cutoff convergence at
$\Delta v=2.5\,\mathrm{m/s}$.}
\label{tab:velocity-cutoff}
\small
\begin{tabular}{c r r r r}
\toprule
Cutoff & $G_p$ & $g_{c,{\rm hist}}$ & $\Delta_{\rm hist}$ &
Relative $G_p$ error\\
\midrule
$3.0\sigma_v$ & 19.36195 & 19.46857 & $-0.10662$ & $2.248\%$\\
$3.5\sigma_v$ & 19.11527 & 19.21041 & $-0.09514$ & $0.945\%$\\
$4.0\sigma_v$ & 18.98214 & 19.07172 & $-0.08957$ & $0.242\%$\\
$4.5\sigma_v$ & 18.93642 & 19.02414 & $-0.08773$ & $6.4\times10^{-4}\%$\\
$5.0\sigma_v$ & 18.93629 & 19.02402 & $-0.08773$ & reference\\
\bottomrule
\end{tabular}
\end{table}

For the experimental phase-matching angle, the separate slow reference now
integrates Eq.~\eqref{eq:two-dimensional-velocity} and passes the grid/cutoff
criteria in Table~\ref{tab:angular-doppler-convergence}. Its angular rms width
$\sigma_{\delta,\theta}/(2\pi)=1.38017\MHz$ is not negligible relative to
$\gamma_{gg}/(2\pi)=0.103\MHz$ and, as Table~\ref{tab:pole-spectrum} shows,
is comparable to the power-broadened Raman pole width itself. This result does
not certify the one-dimensional finite-seed production scan.

\clearpage
\section{Commutator-preservation test}

At the literature operating point, assume for this diagnostic that
the collected probe and conjugate have equal effective Gaussian mode areas,
$A_p=A_c=1.7105972\times10^{-7}\,\mathrm{m^2}$. Their distinct optical
frequencies are retained in $Q$, rather than replaced by $Q\propto I$.
The resulting photon-flux-normalized static drift in the
symmetric mismatch gauge is
\begin{equation}
M_q=
\begin{pmatrix}
-1.795444+8.518638\ii&-0.593347+297.229796\ii\\
0.962238-297.200894\ii&0.508740+60.465455\ii
\end{pmatrix}\ {\rm m}^{-1}.
\label{eq:corrected-literature-drift}
\end{equation}
If the actual conjugate collected-mode area differs from the equal-area audit
assumption, $Q^{-1}M_EQ$ must be recomputed before applying the canonical
metric $J$. The normalization operation itself is now implemented; the
equal-area apparatus input remains conditional.
For the numerical test only, factor the Hermitian matrix
$K=-(M_qJ+JM_q^\dagger)$ as $K=BJ_FB^\dagger$. Here its eigenvalues are
$0.285202\,{\rm m}^{-1}$ and $4.323166\,{\rm m}^{-1}$. Both are positive, so
by Sec.~\ref{sec:realizability} the displayed independent-vacuum algebraic
completion with
$J_F=I$ exists: at this point the drift requires loss-type reservoir
channels.
This construction proves that the drift can be embedded in a canonical
channel; it is not a microscopic derivation of atomic noise.

\begin{table}[ht]
\centering
\caption{Canonical-commutator audit.}
\label{tab:commutator-test}
\small
\begin{tabularx}{\textwidth}{L{4.2cm}Y L{2.0cm}}
\toprule
Test & Residual & Result\\
\midrule
Option-A drift with reservoir omitted &
$\max|TJT^\dagger-J|=3.388$ & \fail\\
Local identity with the explicit $Q$ transform &
$\max|M_qJ+JM_q^\dagger+BJ_FB^\dagger|=4.44\times10^{-16}\,{\rm m}^{-1}$ &
\pass\ conditional\\
Integrated identity through the cell, explicit $Q$ &
$\|TJT^\dagger+\int T_zBJ_FB^\dagger T_z^\dagger\dd z-J\|_F/
\|J\|_F=3.71\times10^{-13}$ & \pass\ conditional\\
Microscopic $M_q(\Omega)$, $B(\Omega)$, and $D(\Omega)$ over 0.1--4 MHz &
Not implemented in the current four-level code & \pending\\
\bottomrule
\end{tabularx}
\end{table}

Because canonical commutator preservation constrains only the antisymmetric
reservoir correlations, it fixes neither the symmetrized diffusion matrix
$D(\Omega)$, nor its temperature, nor excess spontaneous-emission noise; the
vacuum-fixture values below therefore serve as illustrative diagnostics for one
chosen vacuum completion, not as bounds or Langevin-corrected physical
predictions.

\section{Legacy gain-only diagnostic and its failure domain}
\label{sec:gain-referred-ids}

The canonical ideal-amplifier limit must be separated from the historical
gain-only scan score. Define the conjugate photon-flux coefficient
\begin{equation}
g_c^{(\Phi)}\equiv |T_{{\rm cl},21}|^2
=\frac{\omega_p}{\omega_c}G_c .
\label{eq:conjugate-flux-gain}
\end{equation}
For the two-mode map
\[
\hat a_p^{\rm out}=\mu\hat a_p^{\rm in}
+\nu\hat a_c^{{\rm in}\dagger},\qquad
|\mu|^2=G_p,\quad |\nu|^2=g_c^{(\Phi)},
\]
the output commutator is
$G_p-g_c^{(\Phi)}=\mathcal G_{\rm flux}$. Only when
$\mathcal G_{\rm flux}=1$ is this map canonical without additional reservoir
operators. In that lossless limit the photon-number difference is conserved,
and a bright coherent probe seed with a vacuum conjugate input gives
\begin{equation}
\boxed{\;
S_{\rm ideal}=\frac{1}{G_p+g_c^{(\Phi)}}
=\frac{1}{2G_p-1}.\;}
\label{eq:ids-ideal}
\end{equation}
Symmetric external efficiency $\eta$ then gives
\begin{equation}
S(\eta)=1-\eta+\eta S_{\rm ideal}.
\label{eq:ids-symmetric-loss}
\end{equation}

For an ideal canonical amplifier with unequal efficiencies or a weighted
difference, set $A=G_p$, $B=g_c^{(\Phi)}$, and $A-B=1$. In units of the incident
seed photon flux, the exact bright-seed variance and SQL are
\begin{align}
\mathcal V_w={}&\eta_p^2A(A+B)+\eta_p(1-\eta_p)A\nonumber\\
&+w^2\!\left[\eta_c^2B(A+B)+\eta_c(1-\eta_c)B\right]
-4w\eta_p\eta_cAB,\\
S_w={}&\frac{\mathcal V_w}{\eta_pA+w^2\eta_cB} .
\label{eq:ids-balanced}
\end{align}
Here $w=1$ is a raw difference and $w=\eta_pA/(\eta_cB)$ balances the detected
DC levels. For unequal losses the weight that minimizes normalized noise is
covariance dependent; it cannot in general be inferred from detected powers
alone. Equation~\eqref{eq:ids-balanced} reduces to
Eq.~\eqref{eq:ids-symmetric-loss} when $w=1$ and
$\eta_p=\eta_c$.

The archived optimizer instead evaluated the numerical score
\begin{equation}
S_{\rm legacy}=\operatorname{clip}_{[0,1]}
\frac{(G_p-g_c^{(\Phi)})^2}{G_p+g_c^{(\Phi)}} .
\label{eq:ids-legacy}
\end{equation}
This is a definition retained only to reproduce historical scans. It is not a
covariance and its clipping does not prevent passive points from giving
$S_{\rm legacy}<1$; for example $(G_p,g_c^{(\Phi)})=(0.2,0.1)$ gives
$S_{\rm legacy}=0.033$. For the dissipative cell,
$\mathcal G_{\rm flux}\ne1$ requires the reservoir channels of
Secs.~\ref{sec:realizability} and \ref{sec:lyapunov}. Unequal detector weights
must then be applied to the full covariance through Eq.~\eqref{eq:spectrum},
not reconstructed from the two mean gains.

\clearpage
\section{Model point and experimental comparison}

This benchmark applies the passive-limit-tested Option-A propagation signs and
$N_F=3$ to the inherited finite-reference-field four-level Floquet response.
It is not an output of the pump-null-space and frequency-resolvent formulation
of Parts~II--III. It uses
a $5\sigma_v$ one-dimensional velocity range and the documented point
\[
\frac{\Delta}{2\pi}=+0.90\GHz,\quad
\frac{\delta}{2\pi}=-8\MHz,\quad T=121\degC,\quad L=12.5\,{\rm mm},
\]
with $600\,{\rm mW}$ pump power, $8\,\mu{\rm W}$ probe-seed power,
$w_{\rm pump}=530\,\mu{\rm m}$, $w_p=330\,\mu{\rm m}$, and
$\theta=0.32^\circ$. The radii are $1/\e^2$ intensity radii.
The calculation omits the transverse two-photon Doppler average and
therefore remains a reduced-model benchmark. In the same symmetric mismatch
gauge used for Eq.~\eqref{eq:corrected-literature-drift}, the transfer is
\[
T_E=
\begin{pmatrix}
18.839694+6.971473\ii&-8.603522+18.218019\ii\\
8.625460-18.205907\ii&17.483771+10.224863\ii
\end{pmatrix}.
\]
For equal collected mode areas, this gives
\[
G_p=403.53550,\quad G_c=405.85359,\quad
g_c^{(\Phi)}=405.84704,\quad \mathcal G_{\rm flux}=-2.31154.
\]
Here $T_E$ is the electric-field transfer; applying the explicitly constructed
$Q^{-1}T_EQ$ gives $g_c^{(\Phi)}=(\omega_p/\omega_c)G_c$ with
$\omega_c/\omega_p=1.00001614$. The former $(32.56066,31.23207)$ result used
the opposite susceptibility-propagation sign and fails the passive
Beer--Lambert test. The $(5.61434,4.66816)$ result used the passive
susceptibility sign but retained both the duplicated
$p_F/[2(2I+1)]$ population factor and the wrong symmetric-frame mismatch
sign. The intermediate $(107.64682,107.36704)$ result removed the population
duplicate but still used that wrong mismatch sign. The later
$(1054.30596,1061.88534)$ result corrected those conventions but represented
the inherited transit floor as coherence-only damping. The current result adds
the trace-preserving thermal reload channel. None of the earlier numbers is
retained after the combined carrier, mismatch, normalization, and transit audit.
The inherited residual $\ell_s=0.74$ has not been refitted to compensate.

\begin{table}[ht]
\centering
\caption{Separation of mean-field gain and quantum-noise layers at the
literature operating point, using explicit $Q$ with an equal collected-area
assumption. The symmetric external efficiency used in the diagnostic
columns is $\eta_{\rm ext}=0.8694$.}
\label{tab:model-layer-comparison}
\small
\begin{tabularx}{\textwidth}{L{3.7cm}C{1.35cm}C{1.35cm}C{1.8cm}C{1.8cm}Y}
\toprule
Layer & $G_p$ & $G_c$ & $S_-^{\rm out}$ & $S_-^{\rm det}$ &
Interpretation\\
\midrule
Semiclassical reduced model & 403.535 & 405.854 & --- & --- &
Mean-field gain only\\
Ideal Bogoliubov amplifier matched to $G_p$ & 403.535 & 402.542 &
0.001241 ($-29.06\dB$) & 0.13168 ($-8.80\dB$) &
Lossless limiting case, not the dissipative medium\\
Static independent-vacuum completion fixture & 403.535 & 405.854 &
0.02414 ($-16.17\dB$) & 0.15158 ($-8.19\dB$) &
Canonical diagnostic with $w=1$; $D$ is not microscopic\\
Static independent-vacuum completion fixture, DC balanced & 403.535 & 405.854 &
0.005233 ($-22.81\dB$) & --- &
Same diagnostic with $w\simeq0.99429$; detector rematching not evaluated\\
Microscopic quantum-Langevin model & 403.535 & 405.854 &
\pending & \pending & Required physical calculation\\
Experiment \cite{Sim2025} & $111/8=13.875$ & $109/8=13.625$ & --- &
0.166 ($-7.8\dB$) & Raw $8\,\mu{\rm W}$ input and $111/109\,\mu{\rm W}$
outputs; both power ratios are independently available\\
\bottomrule
\end{tabularx}
\end{table}

The static completion and detector model are flat by construction, so
repeating them at several analysis frequencies adds no information. The
measured features instead show what is missing: technical excess noise below a
few hundred kHz, $-8.7\dB$ from the IDS/SQL slope ratio at $150\kHz$, a
laser-source noise peak near $2\MHz$, and loss of squeezing beyond roughly
$3.5\MHz$ \cite{Sim2025}. None can be reconstructed from two gains, so the
table reports no physical $S_-(\Omega_{\rm SA})$ prediction.

\begin{table}[ht]
\centering
\caption{Model--experiment comparison after correcting propagation signs and
removing duplicate manifold-population weighting.}
\label{tab:discrepancies}
\small
\begin{tabularx}{\textwidth}{L{3.0cm}L{3.5cm}L{3.2cm}Y}
\toprule
Observable & Reduced model & Experiment or benchmark & Discrepancy\\
\midrule
$G_p$ & 403.535 & 13.875 from raw powers & $+2808.4\%$\\
$G_c$ & 405.854 & 13.625 from raw powers & $+2878.7\%$\\
$\delta^*/(2\pi)$ & $-275\MHz$ in the archived reduced scan &
$-8\MHz$ measured operating point & $-267\MHz$; different objective and
incomplete geometry make this nonpredictive\\
$\Delta^*/(2\pi)$ & $-1.50\GHz$ in the archived reduced scan &
$+0.90\GHz$ measured operating point & $-2.40\GHz$\\
Maximum squeezing & no physical prediction & $-7.8\dB$ &
Microscopic diffusion and detector response missing\\
Squeezing bandwidth & no physical prediction; separate Lorentzian pole and
Gaussian Doppler scales are tabulated & approximately $3.5\MHz$ &
Two-dimensional mean-field reference complete; microscopic
$S_-(\Omega_{\rm SA})$ comparison pending (Sec.~\ref{sec:pole-spectrum})\\
\bottomrule
\end{tabularx}
\end{table}

The point
$(\Delta/2\pi,\delta/2\pi,T)=(-1.50\GHz,-280\MHz,110\degC)$ is classified
only as a \emph{numerical optimum under the archived reduced model}. It is
not an experimentally established optimum. Likewise, the numerical proximity
of the chosen-vacuum algebraic fixture to the measured $-7.8\dB$ is not
validation: the former contains no microscopic atomic diffusion or detector
spectrum.

\section{Required limiting cases}

\begin{table}[!ht]
\centering
\caption{Analytic limiting-case checks for the layered model.}
\label{tab:limits}
\small
\begin{tabularx}{\textwidth}{L{3.5cm}Y L{2.2cm}}
\toprule
Limit & Required result & Status\\
\midrule
No atom--field coupling &
$\chi_{ij}\to0$ implies $G_p\to1$, $G_c\to0$, $D\to0$, and $S_-\to1$ &
\pass\ analytically\\
Ideal lossless parametric amplifier &
$M_qJ+JM_q^\dagger=0$, $B=D=0$, and $\mathcal G_{\rm flux}=1$ &
\pass\ analytically\\
Bright probe-seeded ideal amplifier &
$S_{\rm ideal}=1/[G_p+g_c^{(\Phi)}]$, reducing to $1/(2G-1)$ when
$\mathcal G_{\rm flux}=1$ &
\pass\ analytically\\
Symmetric external efficiency &
$S_-^{\rm det}=1-\eta+\eta S_-^{\rm out}$ &
\pass\ analytically\\
Zero phase-matching angle &
$\theta\to0$ removes $k\theta\sigma_v$ from the two-photon Doppler width &
\pass\ analytically\\
Zero dissipation &
The Langevin contribution vanishes and $TJT^\dagger=J$ &
\pass\ as a formal limit\\
Zero cell length &
$W(0)=0$ from Eq.~\eqref{eq:differential-lyapunov}, so
$V_{\rm out}=V_{\rm in}$ &
\pass\ analytically and numerically\\
Undamped two-level drive &
Poles at $0,\pm\ii\sqrt{\Omega^2+\Delta^2}$ from
Eq.~\eqref{eq:torrey-cubic} & \pass\ numerically\\
Vanishing drive in the pole sum &
$\mathcal V_k\to0$ makes every residue $A_m^{(jk)}$ vanish, so
$\chibar\to0$ & \pass\ analytically\\
\bottomrule
\end{tabularx}
\end{table}

\clearpage
\section{Conclusions and validation status}

The definitive status is the hierarchy in Table~\ref{tab:validation-hierarchy}.
The analytic core --- the trace-zero resolvent, the diagonalizable-case pole
sum with its Jordan/contour qualification, constant-matrix transfer,
commutator identities, Lyapunov solution, and external detection map --- is
closed, with the identities in Table~\ref{tab:analytic-verification}
verified to $10^{-12}$ or better. The production finite-Floquet layer now adds an
explicit full-scan adjacent-order gate; the separate velocity audits establish
numerical convergence only for their stated reduced problems. State positivity
and the exact-GKSL half-plane condition remain necessary invariants, not
substitutes for convergence or physical validation.
The separate minus-branch pump reference additionally closes the executable
trace-one state, trace-zero response, finite-seed-limit, and direct pole-sum
checks at the errors reported in Table~\ref{tab:pump-reference-diagnostics}.
It remains a slow reference and does not change the finite-seed production
default; the inherited plus-branch reference remains unsupported.
The separate tensor-Maxwell wrapper additionally closes the declared
two-dimensional grid and cutoff checks for a single crossing angle, with the
lab beat held fixed.

At the literature operating point, correcting the Maxwell propagation signs
changes the reduced-model gains to those of Table~\ref{tab:discrepancies}; their
large disagreement with the independently reported output powers precludes
quantitative gain validation. The independent-vacuum completion is retained
solely as a commutator-preserving algebra test. No gain-only number in this report is
a squeezing prediction, and the separate Raman-pole and angular-Doppler scales
do not determine a squeezing bandwidth.

A physical $S_-(\Omega_{\rm SA})$ still requires several coupled inputs:
promotion of the certified minus-branch atomic response into a production
frequency-dependent drift (and a valid treatment of any claimed plus branch),
microscopic reservoir coupling and diffusion, promotion of the certified
two-dimensional response into the finite-seed/segmented production path plus a
measured beam angular distribution, collected-mode normalization, and measured
detector transfer and electronic-noise functions.
Until those are implemented together, the document is a semianalytic
mean-field reconstruction and quantum-model specification, not a validated
squeezing calculation.

\appendix

\section{Analytically solvable limit: the damped two-level rotor}
\label{sec:torrey}

The full four-level Liouvillian is $16\times16$ and its characteristic
polynomial has degree $16$, so its roots are computed numerically. A reduced
two-level case supplies the physical labels used in
Sec.~\ref{sec:pole-expansion}. With longitudinal and transverse relaxation
times $T_1$ and $T_2$, detuning $\Delta$, and Rabi frequency $\Omega$, the
Bloch generator in the $(u,v,w)$ basis is
\begin{equation}
M_{\rm B}=
\begin{pmatrix}
-1/T_2 & -\Delta & 0\\
\Delta & -1/T_2 & \Omega\\
0 & -\Omega & -1/T_1
\end{pmatrix},
\label{eq:bloch-generator}
\end{equation}
with the factored characteristic polynomial
\begin{equation}
\boxed{\;
\left(\lambda+\tfrac1{T_2}\right)
\left[\left(\lambda+\tfrac1{T_2}\right)\left(\lambda+\tfrac1{T_1}\right)
+\Omega^2\right]
+\Delta^2\left(\lambda+\tfrac1{T_1}\right)=0 .}
\label{eq:torrey-cubic}
\end{equation}
Its informative limits are immediate:
\begin{itemize}
\item \textbf{Undriven}, $\Omega=0$: the population pole is $-1/T_1$ and the
free-precession pair is $-1/T_2\pm\ii\Delta$.
\item \textbf{Resonant}, $\Delta=0$: the $u$ quadrature has pole $-1/T_2$;
the remaining quadratic gives a damped-Rabi pair only when
$|\Omega|>|T_1^{-1}-T_2^{-1}|/2$, and otherwise gives two additional real
overdamped branches.
\item \textbf{Undamped}, $T_1,T_2\to\infty$:
\begin{equation}
\lambda_0=0,\qquad
\lambda_\pm=\pm\ii\Omega',\qquad
\Omega'=\sqrt{\Omega^2+\Delta^2}.
\label{eq:generalized-rabi}
\end{equation}
\end{itemize}
Thus the roots comprise one population-recovery branch and either a
complex-conjugate damped-Rabi pair or two additional real overdamped branches.
This is the structure named by Eq.~\eqref{eq:pole-parametrization} in the
multilevel case. Equation~\eqref{eq:torrey-cubic} follows Torrey's classical
transient-nutation analysis \cite{Torrey1949}.

\section{Dimensional-consistency audit}

\begin{table}[ht]
\centering
\caption{Dimensional audit in SI units.}
\label{tab:dimension-audit}
\small
\begin{tabularx}{\textwidth}{L{4.1cm}L{3.0cm}Y}
\toprule
Quantity & Dimension & Check\\
\midrule
$\Delta,\delta,\Gamma,\Omega_{\rm SA},\Delta_e,\lambda_m$ &
$\mathrm{s}^{-1}$ & Angular frequencies and Liouvillian poles\\
$\chibar_{ij}$; $Nd^2/(\epsilon_0\hbar)$ &
$\mathrm{s}$; $\mathrm{s}^{-1}$ & Their product is dimensionless $\chi_{ij}$\\
$M_{\rm cl},M_q,D,K$ & $\mathrm{m}^{-1}$ & Spatial drift, diffusion, and
commutator defect per unit length\\
$\mu_\pm$ & $\mathrm{m}^{-1}$ & Eigenvalues of a spatial generator\\
$W(L)$ & $1$ & $\int D\,\dd z$ with dimensionless covariance\\
$V(x;\sigma,\gamma)$ & $\mathrm{s}$ & $-\pi V$ is an inverse angular frequency\\
$G_p,G_c,S_-$ & $1$ & Power gains and normalized noise spectrum\\
\bottomrule
\end{tabularx}
\end{table}

\section{Assumptions and validity domains}
\label{app:assumptions}

\begin{longtable}{L{3.4cm}L{5.0cm}L{5.6cm}}
\caption{Assumption and validity-domain ledger.}
\label{tab:assumptions}\\
\toprule
Layer & Assumption & Validity domain and failure mode\\
\midrule
\endfirsthead
\toprule
Layer & Assumption & Validity domain and failure mode\\
\midrule
\endhead
\midrule
\multicolumn{3}{r}{\small Continued on next page}\\
\endfoot
\bottomrule
\endlastfoot
Atomic basis & Four hyperfine manifolds with Zeeman structure averaged into
effective strengths & Hyperfine-scale qualitative trends; fails for
polarization, optical pumping, dark states, and magnetic response\\
Pump-only weak-response reference & Undepleted prescribed pump, static
four-level Markovian Liouvillian, trace-one null space, and trace-zero Nambu
resolvent & Executable slow reference for the minus Raman branch only; not used
by the current gain tables, plus branch unsupported, and controlled only under
the stated phenomenological decay channels; a separate 2-D wrapper is available\\
Current gain-state path & Trace-one finite-seed density matrix; production
$N_F\geq3$ with an $N_F-1$ full-scan gate, audit $N_F=1,2,3$ & Not a pump-only state; convergence and finite-seed
back-action must be reported separately\\
Historical rate/Sylvester model & Populations solved without coherence feedback &
Weak saturation or rapid dephasing; not an exact strong-pump steady state\\
Weak fields & Probe and conjugate enter to first order & Valid below the onset
of pump depletion, gain saturation, and nonlinear probe back-action\\
Pole expansion & $\mathcal L$ is diagonalizable & Exact where it holds;
defective Liouvillian eigenvalues require a Jordan or contour form with
higher-order poles, and eigen-residues become ill conditioned near defective
or nearly defective points\\
Lorentzian reading of poles & Residues vary slowly over the analysis band &
Overlapping poles with comparable residues interfere and do not decompose
into independent visible peaks\\
Floquet expansion & Finite $N_F$ & Must be increased until observables
converge; $N_F=1$ fails at the archived point, and no positivity test
detects the failure\\
Velocity average & Maxwell distribution, independent Cartesian components &
The separate minus-branch reference resolves $(v_z,v_x)$ for one declared
angle; the finite-seed production average remains one-dimensional and omits
beam divergence\\
Classical propagation & $z$-independent coefficients and fixed density,
temperature, and pump & Closed-form exponential; fails with longitudinal
pump attenuation, density gradients, or depletion\\
Lyapunov solution & Constant $M_q$, $D$ over the integration interval &
No condition on $M_q$; use the segmented recursion for varying coefficients.
Only the eigenbasis \emph{resolution} additionally needs a well-conditioned
$S$\\
Phase convention & Diagonal complex susceptibility plus vacuum/geometric
mismatch only & Avoids double counting of $\Rea\chi$\\
Quantum propagation & Linear Gaussian Langevin channel & Requires a
microscopic reservoir matrix and appropriate Markov approximation\\
Static vacuum completion & Minimal-port independent-vacuum factorization of
the drift & Commutator fixture only; neither a universal Caves-bound claim nor
an atomic-noise prediction, and frequency independent by construction\\
Detection & External path and detector losses act after the cell &
Valid only for losses physically distinct from distributed atomic absorption\\
Electronics & Linear transfer functions and calibrated SQL reference &
Requires measured $H_p,H_c,w,S_{\rm el}$ at every analysis frequency\\
\end{longtable}

\section{Parameter provenance and sensitivity}
\label{app:parameters}

The classifications below distinguish quantities measured for the cited
apparatus, assumed in the reduced model, fitted or inherited
phenomenologically, and algebraically derived. A dash means that the
parameter does not enter that model layer. Unless an experimental uncertainty
and its source are stated explicitly, a numerical ``audit range'' below is only
an illustrative sensitivity sweep; it is not a confidence interval, prior, or
measurement uncertainty.

\begingroup
\small
\setlength{\tabcolsep}{3pt}
\begin{longtable}{L{2.8cm}L{3.4cm}L{2.5cm}L{3.4cm}L{3.4cm}}
\caption{Measured, assumed, fitted, inherited, and derived parameters.}
\label{tab:parameter-provenance}\\
\toprule
Parameter & Physical definition & Class and source &
Procedure or allowed range & Sensitivity of $G_p,G_c,S_-$\\
\midrule
\endfirsthead
\toprule
Parameter & Physical definition & Class and source &
Procedure or allowed range & Sensitivity of $G_p,G_c,S_-$\\
\midrule
\endhead
\midrule
\multicolumn{5}{r}{\small Continued on next page}\\
\endfoot
\bottomrule
\endlastfoot
$L=12.5\,{\rm mm}$ &
Vapor interaction length &
Measured; experimental cell \cite{Sim2025} &
Use manufacturing/calibration uncertainty when available &
Gains depend through $\exp(M_{\rm cl}L)$; noise through $W(L)$ of
Eq.~\eqref{eq:lyapunov-closed-form}\\
$T=121\degC$ &
Cell-reservoir operating temperature &
Experimentally measured operating point \cite{Sim2025} &
Not a fitted model optimum &
Strong density and Doppler dependence; must be rescanned\\
$P_{\rm pump}=600\,{\rm mW}$, $P_0=8\,\mu{\rm W}$ &
Incident pump and probe-seed powers used for squeezing &
Experimentally measured operating point \cite{Sim2025} &
The probe-seed was approximately $10\,\mu{\rm W}$ in the apparatus description;
$8\,\mu{\rm W}$ for the cited squeezing trace, with raw outputs
$P_p=111\,\mu{\rm W}$ and $P_c=109\,\mu{\rm W}$ &
Pump changes the steady state and power broadens the poles
(Table~\ref{tab:pole-spectrum}); the probe-seed sets the bright-beam
linearization and the SQL\\
$w_{\rm pump}=530\,\mu{\rm m}$, $w_p=330\,\mu{\rm m}$ &
$1/\e^2$ intensity radii, not diameters &
Measured \cite{Sim2025} &
Use the reported beam convention &
Set Rabi frequency and spatial overlap\\
$\theta=0.32^\circ$ &
Phase-matching angle &
Measured geometry \cite{Sim2025} &
Zero-angle check required &
Produces $1.38\MHz$ rms angular Raman broadening at $121\degC$, comparable to the
power-broadened Raman pole width\\
$\Delta/(2\pi)=+0.90\GHz$,
$\delta/(2\pi)=-8\MHz$ &
Experimental optical and two-photon operating detunings &
Experimentally measured operating point \cite{Sim2025} &
Not identified with the archived numerical optimum &
Corrected-sign, one-use-normalized reduced model gives
$G_p=403.535$, $G_c=405.854$ after the thermal transit reset; the $0.74$
residual is not refitted\\
$\Gamma/(2\pi)=5.746\MHz$ &
Natural optical linewidth used in the model &
Literature benchmark &
Ordinary-frequency value converted once to angular rate &
Sets optical resolvents and Voigt width\\
$\gamma_{\rm opt}/(2\pi)=0.742\MHz$ at $121\degC$ &
Added collisional optical half-width $\tfrac12\beta N$ &
Derived from the density and the self-broadening coefficient &
Scales linearly with $N(T)$ &
Broadens the optical poles of Table~\ref{tab:pole-spectrum}\\
$\sigma_v=193.7\,{\rm m/s}$ at $110\degC$;
$196.5\,{\rm m/s}$ at $121\degC$ &
One-axis thermal rms speed &
Derived from Maxwell statistics &
$\sigma_v=\sqrt{k_BT/m}$ &
Sets both one-photon and mismatch-projected Doppler widths
(one-photon rms $243.6\MHz$ and $247.2\MHz$; angular-mismatch rms
$1.36\MHz$ and $1.38\MHz$, at $110\degC$ and $121\degC$, respectively)\\
$\ell_s=0.74$ &
Effective line-strength or overlap multiplier &
Inherited from code; phenomenological coefficient with unknown fit provenance &
Original fitting record unavailable; illustrative $\pm10\%$ sweep
0.666--0.814, not an uncertainty interval &
At fixed $N_F=3$ atomic response, $(G_p,G_c)$ changes from
$(188.447,189.054)$ through $(403.535,405.854)$ to $(857.836,863.691)$;
physical $S_-$ remains unavailable\\
$\kappa=0.1$ &
Coupling of absorbed pump fraction to detected pump-scattering noise &
Inherited from code; assumed phenomenological coefficient &
No documented fit; illustrative sweep 0--0.2, not an uncertainty interval,
pending apparatus calibration &
No effect on $G_p,G_c$; the diagnostic obeys
$N_{\rm ps}=\kappa(1-\exp[-0.149609])$ with slope $0.138956$, but normalized
$S_-(\Omega)$ requires $F_{\rm ps}$ and detector calibration\\
$\gamma_t/(2\pi)=100\kHz$ at $121\degC$
($\gamma_{gg,\mathrm{total}}/(2\pi)=102.882\kHz$) &
Thermal transit-reset input; the total ground-coherence rate also contains a
$2.882\kHz$ collisional pure-dephasing term &
Inherited from code; assumed phenomenological rate &
Illustrative $\gamma_t/(2\pi)$ sweep 90--110 kHz, not an uncertainty interval;
independent linewidth fit absent &
Recomputed $N_F=3$ endpoints give $(G_p,G_c)=(440.250,442.733)$ at
$\gamma_t/(2\pi)=90\kHz$, $(403.535,405.854)$ at $100\kHz$, and
$(370.626,372.790)$ at $110\kHz$ (total rates $92.882$, $102.882$, and
$112.882\kHz$, respectively);
the pole width is not a direct measurement of this parameter\\
$1-\eta_{\rm path}=5.5(2)\%$ &
Post-cell optical-path loss &
Measured \cite{Sim2025} &
Allowed range 5.3--5.7\%; not cell absorption &
No effect on gains; for symmetric loss
$\partial S_{\rm det}/\partial\eta=S_{\rm out}-1$\\
$1-\eta_{\rm det}=8.0(2)\%$ &
Detector and remaining system inefficiency &
Measured \cite{Sim2025} &
Allowed range 7.8--8.2\% &
No effect on gains; combines multiplicatively with path efficiency\\
$\eta_{\rm ext}=0.8694$ &
Illustrative total external efficiency
$0.945\times0.920$ &
Derived from measured losses &
The paper also quotes total loss $13.5(2)\%$; definitions must be reconciled &
Maps $S_{\rm out}$ to $1-\eta_{\rm ext}+\eta_{\rm ext}S_{\rm out}$\\
$N_F$ &
Retained Floquet sideband order &
Numerical truncation &
$N_F=3$ production default with a full-scan $N_F=2$ comparison &
$N_F=1$ changes the historical coefficients by 13--17\% and reverses the sign
of their gap, without
violating any positivity test\\
$\ell_{\rm eff}=0.74/12$ &
Residual line/overlap factor times the sublevel average &
Population supplied by trace-normalized $\rho_{\rm ss}$; duplicate removed &
$\ell_{\rm eff}=0.0616667$; see Appendix~\ref{app:normalization} &
Representative-state coefficient test matches the AutoOD structural factor\\
\end{longtable}
\endgroup

No entry involving $\ell_s$, $\kappa$, or the adopted $\gamma_{gg}$ is
described as exact. Their source records and any apparatus-level fitting
provenance remain required metadata.
The atomic solve also begins from an already-declared cell-input probe power;
it contains no EOM $V_\pi$, insertion-loss, sideband-fraction, or impurity
calibration. Missing EOM provenance therefore blocks an apparatus-level mapping
from RF settings to that optical input and blocks any EOM-noise claim.

These missing records activate the claim gate rather than being filled by an
illustrative sweep: algebraic identities and finite-model convergence may still
be reported, but quantitative apparatus-level gain remains
\textsc{unsupported} and physical $S_-(\Omega_{\rm SA})$ remains
\textsc{unavailable}. A calibration-point agreement, if later obtained, may not
promote either claim without a predeclared held-out comparison.

\section{Dipole and population normalization ledger}
\label{app:normalization}

Let $d_{\rm D1}=\langle J_e\|er\|J_g\rangle$ be the fine-structure reduced
matrix element. The archived coefficient is not the polarization-summed
hyperfine strength. It is the following \emph{AutoOD $q=0$ convention}:
\begin{align}
\mathcal S_0(F_g,F_e)
&=\sum_{m=-\min(F_g,F_e)}^{+\min(F_g,F_e)}
\begin{pmatrix}F_e&1&F_g\\m&0&-m\end{pmatrix}^{\!2},\\
C_F^2&=(2F_g+1)(2F_e+1)(2J_g+1)(2J_e+1)(2L_g+1)
\mathcal S_0(F_g,F_e)\nonumber\\
&\quad\times
\begin{Bmatrix}J_g&J_e&1\\F_e&F_g&I\end{Bmatrix}^{\!2}
\begin{Bmatrix}L_g&L_e&1\\J_e&J_g&S\end{Bmatrix}^{\!2}.
\label{eq:CF-normalization}
\end{align}
For $I=5/2$, $S=1/2$, $L_g=0$, $L_e=1$, and
$J_g=J_e=1/2$. Then $\mathcal S_0=1/3$, which gives the stored exact values
\[
C_F^2=\{10,35,35,28\}/81
\]
for $(F_g,F_e)=(2,2),(2,3),(3,2),(3,3)$, respectively.

The AutoOD orbital reduced dipole and the fine-structure reduced dipole used
by GABES obey $|d_{\rm orb}|^2=3|d_{\rm D1}|^2$ in this D1 convention.
Therefore
\[
|d_{\rm eff}|^2=C_F^2|d_{\rm orb}|^2
=3C_F^2|d_{\rm D1}|^2,
\]
and the dimensionless amplitude multiplier that accompanies
$d_{\rm D1}$ is
\begin{equation}
a_{ge}=\sqrt{3C_F^2}.
\label{eq:age-normalization}
\end{equation}
This derivation explains the code factor but also its limitation:
Eq.~\eqref{eq:CF-normalization} is a $q=0$, Zeeman-averaged absorption
convention, not a universal incoherent sum over $q=0,\pm1$. A
polarization-resolved Liouvillian must replace it with explicit
Clebsch--Gordan amplitudes and must not multiply by $a_{ge}$ again.

The formerly archived effective factor was assembled as
\begin{equation}
\ell_{\rm eff}^{\rm old}
=\ell_s\,\frac{p_{F=3}}{2(2I+1)}
=0.74\times\frac{7}{144}
=0.0359722 .
\label{eq:ell-effective}
\end{equation}
This double-counted the manifold population: the trace-normalized
$\rho_{\rm ss}$ already supplies the actual, pump-modified population of the
representative $F=2,3$ states. The corrected one-use normalization is
\begin{equation}
\boxed{\ell_{\rm eff}=\ell_s\frac{1}{2(2I+1)}
=\frac{0.74}{12}=0.0616667 .}
\label{eq:ell-effective-corrected}
\end{equation}
In a thermal reference state the response itself contributes
$\rho_{FF}=p_F$, yielding the AutoOD structural coefficient
$3C_F^2p_F/[2(2I+1)]$ without a second external $p_F$. The residual $\ell_s$
remains phenomenological and must not be silently refitted to restore an old
gain value.

The normalization rule for future implementations is:
\begin{enumerate}
\item use the fine-structure reduced dipole matrix element once;
\item use either explicit Clebsch--Gordan matrix elements or the averaged
$a_{ge}$, never both;
\item include isotope abundance once in $N$;
\item obtain ground-manifold population either from $\rho_{\rm ss}$ or an
external $p_F$, never both;
\item include spontaneous-emission branching ratios in the collapse
operators, not again in susceptibility amplitudes; and
\item apply external optical efficiencies only in the detection layer.
\end{enumerate}

\begin{table}[ht]
\centering
\caption{Code-path normalization audit for the current reduced FWM solver.}
\label{tab:normalization-audit}
\small
\begin{tabularx}{\textwidth}{L{3.2cm}Y L{2.5cm}}
\toprule
Factor & Where it enters & Verdict\\
\midrule
Reduced D1 dipole $d_{\rm D1}$ &
Once in the macroscopic conversion in
\texttt{gabes/observables.py} & no duplicate found\\
Hyperfine/Zeeman strength &
$\sqrt{3C_F^2}$ in the weak-field drive and once in polarization readout;
their product gives $3C_F^2$. It is not also present in
\texttt{physical\_coupling\_norm}. & no algebraic duplicate found\\
Isotope abundance &
The FWM path uses the pure-$^{85}$Rb density from
\texttt{hyperfine.number\_density}; it does not multiply the natural
$0.7217$ abundance. & no duplicate in this path\\
Ground-manifold population &
$\rho_{\rm ss}$ is trace normalized over the two lumped ground manifolds and
supplies the population once. \texttt{physical\_coupling\_norm} now supplies
only $1/[2(2I+1)]$. & algebraic duplicate removed; structural test\\
Spontaneous branching &
The $C_F^2$-normalized decay fractions occur in Liouvillian collapse
channels; no additional branching multiplier is applied in
\texttt{chi\_phys}. & no duplicate found\\
Residual $\ell_s$ &
Multiplies the macroscopic coupling once. & inherited, not derived\\
\bottomrule
\end{tabularx}
\end{table}

The coefficient ledger now removes the algebraic population duplication and a
representative thermal-state/AutoOD coefficient identity is regression-tested.
Executable tests now also inspect the pump-state $\rho_{\rm ss}$ and its
trace-zero weak response. Those dynamical tests do not supply an independent
absolute-OD scale and therefore are not, by themselves, a full weak-probe OD
validation. Any future end-to-end OD comparison must retain the same one-use
population-counting rule.

\section{Four-level model limitations and unresolved physics}
\label{app:limitations}

The four-hyperfine-level model is suitable for qualitative gain structure,
reduced hyperfine-scale intuition, and comparison with the existing reduced
simulation. It is not automatically sufficient for absolute squeezing,
polarization-sensitive effects, magnetic-field response, quantitative
spontaneous-emission noise, or accurate Raman linewidths.

Physics outside the current basis and reservoir model includes:
\begin{itemize}
\item polarization-dependent Clebsch--Gordan pathways and explicit
$\sigma^\pm$ and $\pi$ selection rules;
\item Zeeman optical pumping, ground-state orientation and alignment,
dark-state formation, and magnetic-field-dependent coherences;
\item interference among distinct $m_F$ pathways and
polarization-dependent spontaneous emission;
\item spatial pump depletion, longitudinal attenuation, transverse
Gaussian averaging, diffraction, mode mismatch, and cell-window etalons;
\item buffer-gas or wall-collision dynamics, transit-time distributions,
radiation trapping, and non-Markovian reservoirs;
\item multidimensional velocity-changing collisions, beam angular
distributions, and promotion of the separate non-collinear Raman-Doppler
reference into the finite-seed segmented production path;
\item a microscopic frequency-dependent atomic diffusion matrix, excess
fluorescence, and atom--reservoir correlations;
\item detector transfer functions, electronic cross-talk, common-mode
rejection, and the apparatus-specific $2\MHz$ laser-source noise peak; and
\item an independently measured SQL calibration and gain definition
matched exactly to the theoretical modes.
\end{itemize}

None of these can be assumed small. The gain discrepancy of
Table~\ref{tab:discrepancies} and the $1.36$--$1.38\MHz$ angular Doppler
scale, which is comparable to the power-broadened Raman pole width itself,
are order-unity effects; the present calculation does not yet define a
squeezing-bandwidth prediction.

\section{Implementation notes for the analytic layers}
\label{app:implementation}

The following are the practical requirements implied by
Secs.~\ref{sec:pole-expansion} and \ref{sec:lyapunov}.

\begin{enumerate}
\item \textbf{Do not diagonalize $M_q$ without a guard.} Use
Eq.~\eqref{eq:van-loan-result} by default. If
Eq.~\eqref{eq:lyapunov-eigenbasis} is used for interpretation, check
$\operatorname{cond}(S)$ and fall back when it exceeds a threshold; the two
forms agree to $10^{-15}$ in the well-conditioned regime.
\item \textbf{Assert the invariants.} $W$ must be Hermitian and positive
semidefinite for any Hermitian, positive-semidefinite $D$, and $W(0)=0$. These
are cheap tests that catch sign and conjugation errors in the
creation/annihilation bookkeeping.
\item \textbf{Treat the stationary subspace explicitly at zero frequency.}
For $\Omega\ne0$ the unique trace-mode residue vanishes by
Eq.~\eqref{eq:zero-residue}. At $\Omega=0$, solve on the trace-zero subspace or
use the Drazin inverse; do not evaluate a formal $0/0$ term. If more than one
stationary mode exists, project the full stationary subspace and check every
zero-mode residue.
\item \textbf{Track residues, not only eigenvalues.} A narrow pole with a
negligible residue is not a visible spectral feature. Ranking poles by
$|A_m|/\Gamma_m$ rather than by $\Gamma_m$ alone identifies which ones
actually shape $S_-(\Omega_{\rm SA})$.
\item \textbf{Recompute $K$ after any normalization change.}
Equation~\eqref{eq:K-matrix} is contracted with $J$ and is therefore
basis dependent through Eq.~\eqref{eq:canonical-rescaling}.
\item \textbf{Report the analysis frequency with every spectrum.} The
$\Omega_{\rm SA}$ dependence enters both $M_q$ and $D$; a spectrum quoted
without its frequency grid cannot be compared with a measurement.
\end{enumerate}

\section{Reproducibility}
\label{app:reproducibility}

The archived Floquet and velocity-grid tables and the separate Option-A
literature-point transfer, Maxwell/$Q$ parity, commutator fixture, and model-layer
comparison, plus the $\ell_s$, $\kappa$, transit-floor, and two-dimensional
Raman-Doppler grid/cutoff sensitivity blocks,
are generated by
\path{analysis/squeezing/analytic_reconstruction/convergence_audit.py}.
Its human-readable and machine-readable outputs are
\path{generated/convergence_audit.md} and
\path{generated/convergence_audit.json}. Those files retain the separate
historical dressed-$k$ scans as numerical-isolation fixtures, while their
literature-point Option-A section uses the current production atomic dissipator
and shared Maxwell/$Q$ constructors. Its Floquet assembly remains analysis-owned
as a cross-check, while production independently supplies general-order
continued-fraction core and compiled-kernel paths. It also reconstructs $M_E$,
$Q$, and $T_E$ independently from literal SI prefactors and
\texttt{scipy.linalg.expm}. At $N_F=3$ the production/reference maximum
differences are $1.64\times10^{-14}$ or smaller. Tables~\ref{tab:floquet-option-a},
\ref{tab:commutator-test}, and \ref{tab:model-layer-comparison} are generated
from that checked path. Passive Beer--Lambert, unrotated-to-symmetric mismatch
equivalence, zero-coupling identity, constant-segment parity, the structural
normalization ledger, an independent stored-$\chi$ Option-A transfer golden,
Bogoliubov $Q$ conversion, and Manley--Rowe energy bounds are pinned in
\path{tests/test_fwm_propagation.py}.
General-order continued-fraction versus dense-block parity, all-harmonic trace
and residual invariants, compiled-kernel parity for $N_F=1,2,3$, and the
full-scan convergence gate are pinned in \path{tests/test_fwm_floquet.py}.
The pump-only state diagnostics, static-frame/Floquet parity, finite-seed limit,
projected-DC response, and direct pole/residue parity are emitted by
\path{convergence_audit.py}; their solver-level fixtures are pinned independently
in \path{tests/test_fwm_pump_response.py}.
The lab-frequency contract, collinear and equal-$k$ limits, analytic rms, and
two-dimensional grid/cutoff checks are pinned in
\path{tests/test_fwm_angular_doppler.py}. The detailed analytic-identity table,
full pole-spectrum table, and physicality-margin table remain report derivations
or separately reported checks and are not described as regenerated by the audit
script.
The document-level status is reported once, in
Table~\ref{tab:validation-hierarchy} and the conclusion, rather than repeated
in a second acceptance table.

\begin{thebibliography}{9}

\bibitem{Sim2025}
G. Sim, H. Kim, and H. S. Moon,
\emph{Intensity-difference squeezing from four-wave mixing in hot
$^{85}$Rb and $^{87}$Rb atoms in single diode laser pumping system},
\emph{Scientific Reports} \textbf{15}, 7727 (2025).

\bibitem{McCormick2007}
C. F. McCormick, V. Boyer, E. Arimondo, and P. D. Lett,
\emph{Strong relative intensity squeezing by four-wave mixing in rubidium
vapor}, \emph{Optics Letters} \textbf{32}, 178--180 (2007).

\bibitem{Turnbull2013}
M. T. Turnbull, P. G. Petrov, C. S. Embrey, A. M. Marino, and V. Boyer,
\emph{Role of the phase-matching condition in nondegenerate four-wave mixing
in hot vapors for the generation of squeezed states of light},
\emph{Physical Review A} \textbf{88}, 033845 (2013).

\bibitem{Lukin1999}
M. D. Lukin, A. B. Matsko, M. Fleischhauer, and M. O. Scully,
\emph{Quantum noise and correlations in resonantly enhanced wave mixing
based on atomic coherence}, \emph{Physical Review Letters} \textbf{82},
1847--1850 (1999).

\bibitem{Caves1982}
C. M. Caves,
\emph{Quantum limits on noise in linear amplifiers},
\emph{Physical Review D} \textbf{26}, 1817--1839 (1982).

\bibitem{Torrey1949}
H. C. Torrey,
\emph{Transient nutations in nuclear magnetic resonance},
\emph{Physical Review} \textbf{76}, 1059--1068 (1949).

\bibitem{VanLoan1978}
C. F. Van Loan,
\emph{Computing integrals involving the matrix exponential},
\emph{IEEE Transactions on Automatic Control} \textbf{23}, 395--404 (1978).

\bibitem{GKS1976}
V. Gorini, A. Kossakowski, and E. C. G. Sudarshan,
\emph{Completely positive dynamical semigroups of $N$-level systems},
\emph{Journal of Mathematical Physics} \textbf{17}, 821--825 (1976).

\bibitem{Faddeeva}
V. N. Faddeeva and N. M. Terent'ev,
\emph{Tables of Values of the Function $w(z)$ for Complex Argument}
(Pergamon, 1961).

\end{thebibliography}

\end{document}
